\documentclass[11pt,a4paper]{article}

% ---------------------------------------------------------------------------
% Packages
% ---------------------------------------------------------------------------
\usepackage[T1]{fontenc}
\usepackage[utf8]{inputenc}
\usepackage{lmodern}
\usepackage[a4paper,margin=2.2cm,headheight=14pt]{geometry}
\usepackage{xcolor}
\usepackage{colortbl}
\usepackage{booktabs}
\usepackage{longtable}
\usepackage{tabularx}
\usepackage{array}
\usepackage{multirow}
\usepackage{enumitem}
\usepackage{hyperref}
\usepackage{titlesec}
\usepackage{fancyhdr}
\usepackage{setspace}
\usepackage{caption}
\usepackage{needspace}
\usepackage{graphicx}
\usepackage{etoolbox}

% ---------------------------------------------------------------------------
% Colours
% ---------------------------------------------------------------------------
\definecolor{navy}{HTML}{1B365D}
\definecolor{steel}{HTML}{2E5984}
\definecolor{p0}{HTML}{1B7A4E}
\definecolor{p1}{HTML}{B45309}
\definecolor{p2}{HTML}{6B7280}
\definecolor{tbd}{HTML}{B91C1C}
\definecolor{v1bg}{HTML}{ECFDF3}
\definecolor{futbg}{HTML}{F3F4F6}
\definecolor{tbdbg}{HTML}{FEF2F2}
\definecolor{rulegray}{HTML}{D1D5DB}
\definecolor{rowalt}{HTML}{F8FAFC}
\definecolor{hdrbg}{HTML}{1B365D}

% ---------------------------------------------------------------------------
% Hyperref / headers
% ---------------------------------------------------------------------------
\hypersetup{
  colorlinks=true,
  linkcolor=steel,
  urlcolor=steel,
  citecolor=steel,
  pdftitle={SoMe App --- User Stories and Requirements Specification},
  pdfauthor={Product Engineering},
  pdfsubject={V1 user stories derived from spec v0.33, meeting transcript, and handwritten mockups}
}

\pagestyle{fancy}
\fancyhf{}
\fancyhead[L]{\small\color{steel}SoMe App --- User Stories \& Requirements}
\fancyhead[R]{\small\color{steel}v1.0}
\fancyfoot[C]{\small\color{steel}\thepage}
\renewcommand{\headrulewidth}{0.4pt}
\renewcommand{\footrulewidth}{0.3pt}

\titleformat{\section}{\Large\bfseries\color{navy}}{\thesection}{0.8em}{}
\titleformat{\subsection}{\large\bfseries\color{steel}}{\thesubsection}{0.7em}{}
\titleformat{\subsubsection}{\normalsize\bfseries\color{navy}}{\thesubsubsection}{0.6em}{}
\titlespacing*{\section}{0pt}{1.6em}{0.7em}
\titlespacing*{\subsection}{0pt}{1.2em}{0.5em}

\setlist[itemize]{leftmargin=1.4em,itemsep=0.2em,topsep=0.3em}
\setlist[enumerate]{leftmargin=1.6em,itemsep=0.25em,topsep=0.3em}
\setstretch{1.08}
\setlength{\parskip}{0.35em}
\setlength{\parindent}{0pt}

\captionsetup{font=small,labelfont=bf}

% ---------------------------------------------------------------------------
% Story helpers
% ---------------------------------------------------------------------------
\newcommand{\priPzero}{\colorbox{p0}{\makebox[4.6em]{\textcolor{white}{\scriptsize\bfseries P0 V1}}}}
\newcommand{\priPone}{\colorbox{p1}{\makebox[4.6em]{\textcolor{white}{\scriptsize\bfseries P1}}}}
\newcommand{\priPtwo}{\colorbox{p2}{\makebox[4.6em]{\textcolor{white}{\scriptsize\bfseries P2 Later}}}}
\newcommand{\priTBD}{\colorbox{tbd}{\makebox[4.6em]{\textcolor{white}{\scriptsize\bfseries TBD}}}}

\newcommand{\srcSpec}{\textit{Spec v0.33}}
\newcommand{\srcTx}{\textit{Transcript}}
\newcommand{\srcMock}{\textit{Mockup 260922}}
\newcommand{\srcAll}{\textit{Spec v0.33; Transcript; Mockup 260922}}

\newenvironment{story}[1]{%
  \Needspace{10\baselineskip}%
  \vspace{0.7em}%
  \noindent
  \colorbox{navy}{\makebox[\dimexpr\textwidth-2\fboxsep\relax][l]{\textcolor{white}{\bfseries\ #1}}}%
  \par\vspace{0.35em}%
}{\par\vspace{0.55em}}

\newcommand{\storymeta}[5]{%
  \noindent
  \begin{tabularx}{\textwidth}{@{}>{\bfseries}l X@{}}
    As a: & #1 \\
    I want to: & #2 \\
    So that: & #3 \\
    Priority: & #4 \\
    Source: & #5 \\
  \end{tabularx}
  \par\vspace{0.3em}
}

\newcommand{\achead}{\noindent\textbf{Acceptance Criteria}\par}
\newcommand{\authz}[1]{\noindent\textbf{Authorization:} #1\par}
\newcommand{\deps}[1]{\noindent\textbf{Dependencies:} #1\par}
\newcommand{\precond}[1]{\noindent\textbf{Preconditions:} #1\par}
\newcommand{\notes}[1]{\noindent\textbf{Notes:} #1\par}

\newcolumntype{L}[1]{>{\raggedright\arraybackslash}p{#1}}
\newcolumntype{Y}{>{\raggedright\arraybackslash}X}

% ---------------------------------------------------------------------------
\begin{document}

% ===========================================================================
% TITLE PAGE
% ===========================================================================
\begin{titlepage}
\centering
\vspace*{1.6cm}
{\color{navy}\rule{\textwidth}{2.2pt}}\\[0.55em]
{\color{steel}\rule{\textwidth}{0.6pt}}\\[1.3em]
{\LARGE\bfseries\color{navy} SoMe App}\\[0.45em]
{\huge\bfseries\color{navy} User Stories and Requirements}\\[0.25em]
{\huge\bfseries\color{navy} Specification}\\[1.0em]
{\color{steel}\rule{0.42\textwidth}{0.8pt}}\\[1.0em]
{\Large Community communication platform}\\[0.35em]
{\large Wire-like chat, community-scoped profiles, verification,}\\[0.15em]
{\large and group tools --- implementation-ready V1 scope}\\[1.6em]
\begin{tabular}{>{\bfseries\color{navy}}l l}
Document type & User Stories / Product Requirements \\
Document ID & SOME-USR-001 \\
Version & 1.0 \\
Status & Draft for team-lead review \\
Date & 22 September 2026 \\
Classification & Internal --- project team \\
\end{tabular}\\[1.4em]
{\color{steel}\rule{\textwidth}{0.6pt}}\\[0.55em]
{\color{navy}\rule{\textwidth}{2.2pt}}\\[1.6em]
\begin{minipage}{0.92\textwidth}
\small
\textbf{Derived from (primary sources of truth)}
\begin{itemize}[leftmargin=1.2em,itemsep=0.15em]
  \item Word specification \textit{Version 0.33 --- 260915}
  \item Team meeting transcript (walkthrough of the specification)
  \item Handwritten UI mockup PDF \textit{SoMe App 260922} (28 screens, IDs AA--D3)
  \item Wire UI screenshots attached to the specification (visual language only)
\end{itemize}
\textbf{This document does not invent product behaviour.} Where sources conflict or are silent, the item is marked \textbf{TBD --- Product Decision Required} or \textbf{Future / Later}.
\end{minipage}
\vfill
{\small\color{p2}Prepared for the development team, product owner, design, and QA.}
\end{titlepage}

% ===========================================================================
% DOCUMENT CONTROL
% ===========================================================================
\section*{Document Control}
\addcontentsline{toc}{section}{Document Control}

\subsection*{Revision history}
\begin{tabularx}{\textwidth}{@{}l l l Y@{}}
\toprule
\textbf{Version} & \textbf{Date} & \textbf{Author} & \textbf{Change} \\
\midrule
1.0 & 2026-09-22 & Product Engineering & Initial user-stories document extracted from spec v0.33, meeting transcript, and handwritten mockup 260922. \\
\bottomrule
\end{tabularx}

\subsection*{Source materials}
\begin{tabularx}{\textwidth}{@{}l l Y@{}}
\toprule
\textbf{ID} & \textbf{Material} & \textbf{How it is used} \\
\midrule
S1 & Spec v0.33 (260915) & Named features, group settings, profile base fields, later flags \\
S2 & Meeting transcript & Intent, V1 vs later, rejected ideas, UX philosophy \\
S3 & Mockup 260922 & Screen inventory, flows, field-level UI, navigation \\
S4 & Wire screenshots & Visual density and interaction patterns only --- not product scope \\
\bottomrule
\end{tabularx}

\subsection*{Implementation context (not product requirements)}
The implementation team currently expects a modern web stack (for example Next.js, TypeScript, Supabase/PostgreSQL, Supabase Auth, Realtime, and Storage). \textbf{Those technologies are not product requirements} unless a source names them. User stories describe behaviour, not vendors.

\subsection*{Reading convention}
\begin{itemize}
  \item \priPzero\ Required for the initial release because the sources place it in the current specification or mockup as current behaviour, and do not mark it later.
  \item \priPone\ Present in sources as current or drawn, but another source delays it, contradicts it, or leaves the exact rule incomplete. Do not implement as if settled.
  \item \priPtwo\ Explicitly later / future / ``not now'' in the sources.
  \item \priTBD\ A product decision is required before the story can be implemented as specified.
\end{itemize}

\tableofcontents
\newpage

% ===========================================================================
\section{Product Overview}
% ===========================================================================

\subsection{One-line product definition}
The product is a mobile-first web application that combines a minimal chat experience (comparable to Wire) with real user profiles, invite-style community isolation, human verification, and form-based discovery, so people in a shared-interest community can talk and meet in real life.

\subsection{Why it exists}
The founding community already uses Wire. Wire provides fast, clean chat but \textbf{has no user profiles}. Members cannot browse who else is in the community, inspect interests, or judge trust from a profile. The intended differentiator, stated in the transcript, is:

\begin{quote}
\textit{Take the Wire app and add user profiles. The rest is quality of life.}
\end{quote}

Quality-of-life additions named in the sources: sticky information, group calendars, easier admin (add/remove members, custom member status), in-app camera tagging of images, private notes/ratings, and community-specific search.

\subsection{Product hierarchy}
\begin{verbatim}
Platform          (the company / whole application --- Platform Admin)
  └── Community   (one interest circle; own URL; own profile-field schema)
        └── Group (a chat room inside that community; own admin and settings)
              └── User (member of the community and of 0+ groups)
                    └── DMs / Posts / Profile / Albums / Notes
\end{verbatim}

A \textbf{community} is not a group. A community has members, groups, and its own profile questions. A \textbf{group} is a chat room inside one community.

\subsection{Design philosophy (source: transcript)}
\begin{itemize}
  \item ``Apple, not Microsoft'': few visible menus; one obvious action per screen.
  \item Most members only post and react.
  \item Power actions live on long-press / full-screen sheets, not dense icon bars.
  \item Dialogues may use most of the screen; do not use tiny hard-to-tap controls.
  \item Visual language follows the attached Wire screenshots (density, avatars, unread badges, tappable chat title) unless a handwritten screen replaces that layout.
\end{itemize}

\subsection{What this product is not (explicitly rejected or later)}
\begin{itemize}
  \item Not a cross-community social network. Search and membership stay inside one community.
  \item Not a free-text global search.
  \item Not Tinder-style kilometre distance search (rejected as too complicated).
  \item Not an automated verification provider or AI identity check in V1.
  \item Not a payments product in V1.
  \item Not a native push-notification product in V1 (browser notifications are later).
\end{itemize}

% ===========================================================================
\section{Actors and Roles}
% ===========================================================================

\begin{tabularx}{\textwidth}{@{}l Y Y@{}}
\toprule
\textbf{Role} & \textbf{Can do (from sources)} & \textbf{Cannot do / limits} \\
\midrule
Platform Admin
  & Everything platform-level. Approve/create communities. For the first community this is the same person as community admin. Can set profile fields for early communities.
  & Sources do not specify a platform console beyond this. \\
Community Admin
  & Verify new users; decide who enters the community; create groups in V1; control community profile fields (personally, for the first communities); community-wide settings as named.
  & Must not access another community's data. Self-serve field UI for other community admins is later. \\
Group Admin
  & Create/rename/delete their group; add/remove members; sticky posts; calendar; group settings; delete any post in that group; custom member status; add a DM counterpart to groups they admin.
  & Cannot change platform/community settings or other groups. Co-admins are later. \\
User
  & Chat in DMs and groups they belong to; edit own profile; post; react; comment; delete own posts; edit own posts if the group allows it; search profiles in their community; leave chats; block in DMs; private notes/ratings; albums.
  & Cannot create groups in V1. Cannot delete others' posts. Cannot see others' private notes/ratings about them. \\
Unapproved user
  & After Join: fill profile and DM the community admin (screen AD). Cannot use the main chat list as an approved member until approved.
  & Exact lock-down of other screens while pending is TBD. \\
\bottomrule
\end{tabularx}

\subsection{User levels 1--3}
The specification states ``user 1--3 levels (community admin or payment level controls the level/rights)''. The only concrete differentiator given in the transcript is \textbf{simple search versus advanced search} (simple = age and gender; advanced = more detail).

\textbf{TBD --- Product Decision Required:} the permission matrix for levels 1--3 beyond that example is not specified. Do not invent entitlements.

\subsection{Authentication versus authorisation}
\begin{itemize}
  \item \textbf{Authentication:} proving who the account is (login / email link).
  \item \textbf{Authorisation:} after login, checking role + scope (this community, this group, this post, this album password).
\end{itemize}
Every write action in V1 requires an authorisation check. Community isolation is the primary privacy boundary.

% ===========================================================================
\section{Scope}
% ===========================================================================

\subsection{In V1 (initial release)}
Login and join; community-scoped routing; pending-approval path; main chat list; DMs and group chats; text and image posts; in-app camera tag; reactions; comments; own delete; own edit if allowed; admin delete-any-post; personal chat notes; group sticky posts; group calendar; V1 group settings; admin-only group creation; add/remove members; custom member status; profiles with base fields plus dynamic field \emph{capability}; form-based community search; albums including lock password; private profile notes and 1--5 star ratings; manual verification.

\subsection{Out of V1 (explicitly later)}
Video/audio/collage posts; disappearing/limited-view posts; browser/OS notifications; co-admins; username-visibility toggle; user-created groups; monetized ads; payments (card, Swish, multi-level fees, profit share); reputation score; profile completeness score; AI-assisted verification; self-serve multi-community onboarding for other admins; Google/social login.

\subsection{Conflicts that block a silent V1 decision}
These items appear in more than one source with incompatible wording. They are \textbf{not} resolved in this document.

\begin{enumerate}[leftmargin=1.6em]
  \item \textbf{Bottom navigation:} transcript says two items (Chats + Search); mockup 260922 draws three (Chats / Profiles / Ads). Own profile is the top-left avatar in both.
  \item \textbf{Who may create groups in V1:} opening spec bullet says users can create a group and be admin; later spec text and the transcript say only the community admin can create a group now.
  \item \textbf{URL model:} spec says ``sub domain'' but the written example is a path (\texttt{https://maindomain/community}; spoken example \texttt{meetirl.com/bikingstockholm}).
  \item \textbf{Account model:} sources never state whether an account is platform-wide (then join communities) or created per community URL.
  \item \textbf{Ads in V1:} transcript says no special ad system --- a normal post, with feed filtering later; mockup draws dedicated screens A6, C4, C9, D1, D2, D3.
  \item \textbf{Direct join:} transcript says not needed at the beginning (always message admin); mockup A9 includes an ``Allow direct join'' checkbox.
  \item \textbf{Screen-ID collisions inside the mockup:} B1 is both Group Calendar and own profile; A3 is both ``your profile'' on A1 and ``other profile'' on the profile page; B5 is both private notes and ``own post menu''.
\end{enumerate}

% ===========================================================================
\section{User Story Methodology}
% ===========================================================================
Stories use: \textit{As a [role], I want [goal], so that [benefit].}

Acceptance criteria are written so QA can turn them into cases. Prefer Given / When / Then. Criteria that cannot be tested as written are marked TBD.

Priority is taken from the sources, not from engineering preference:
\begin{itemize}
  \item Named in spec/mockup as current, not marked later $\rightarrow$ P0.
  \item Drawn or mentioned but contradicted or incomplete $\rightarrow$ P1 or TBD.
  \item Marked later / future / ``not now'' $\rightarrow$ P2.
\end{itemize}

Story IDs are stable. Do not renumber; add new IDs at the end of a prefix.

\subsection{Epic catalogue}
\begin{tabularx}{\textwidth}{@{}l l Y@{}}
\toprule
\textbf{Epic} & \textbf{Prefix} & \textbf{Covers} \\
\midrule
E1 Authentication \& Account Access & AUTH- & Join, login, email link, forgot password, pending state \\
E2 Community Access \& Membership & COMM- & URL, isolation, membership, approval \\
E3 User Profile & PROF- & Own/other profile, albums, images \\
E4 Dynamic Profile Fields & DYN- & Schema, field types, who configures \\
E5 Profile Search & SEARCH- & Form search, filters, results \\
E6 Direct Messaging & DM- & List, start, leave, block, admin shortcut \\
E7 Group Chats & GROUP- & Create, members, settings, leave, join request \\
E8 Posts, Reactions \& Comments & POST- & Feed items in DM and group \\
E9 Verification & VERIFY- & Manual review, in-app camera tag \\
E10 Sticky Posts \& Calendar & STICKY- / CAL- & Group-only extras \\
E11 Personal Notes \& Ratings & NOTE- & Chat notes, profile notes, stars \\
E12 Navigation \& Main UI & NAV- & Shell, +, sorting, philosophy \\
E13 Settings & SET- & Only settings the sources name \\
E14 Notifications & NOTIF- & Later browser notifications \\
E15 Advertising & AD- & V1 concept vs dedicated screens vs later monetisation \\
E16 Payments \& Monetisation & PAY- & Entirely later \\
E17 Administration & ADMIN- & Platform / community / group admin actions \\
E18 Security, Privacy \& Data Access & SEC- & Isolation, GDPR, album gate, private data \\
\bottomrule
\end{tabularx}

\newpage
% ===========================================================================
\section{Epic 1 --- Authentication \& Account Access}
% ===========================================================================
Community landing (AA) offers Login (AB) and Join (AC). Join collects username, password, password confirmation, and email, then sends an email with a login link. Google/social login is explicitly not required now. After login, approved users reach A1; users ``not yet approved'' reach AD.

\begin{story}{AUTH-001 --- Join (create account) from a community URL}
\storymeta{Unregistered visitor}{create an account with username, password, password confirmation, and email}{I can request access to this community}{\priPzero}{\srcMock; \srcSpec}
\precond{The visitor opened a community start page (AA).}
\achead
\begin{enumerate}
  \item Given I am on AA, when I tap Join, then I see AC with username, password, password again, email, and Send.
  \item Given I submit valid values, when Send succeeds, then the system stores a pending account and sends an email containing a login link that opens AB.
  \item Given the passwords do not match or a required field is empty, when I tap Send, then the account is not created and I am shown a validation message.
  \item Given the username or email is already used \emph{in the applicable account scope}, when I tap Send, then the system rejects the join. \textbf{Account scope (per community vs platform-wide) is TBD --- see COMM-006.}
\end{enumerate}
\authz{Unauthenticated. Creates an unapproved account only.}
\deps{COMM-001 (community URL).}
\end{story}

\begin{story}{AUTH-002 --- Email login link after join}
\storymeta{Newly joined user}{receive an email with a login link}{I can open the login screen and continue}{\priPzero}{\srcMock}
\achead
\begin{enumerate}
  \item Given I completed AC, when the email is sent, then it contains a link that opens AB.
  \item The sources do not specify expiry, one-time use, or deep-link token format. \textbf{TBD --- Product Decision Required} for token lifetime and reuse.
\end{enumerate}
\end{story}

\begin{story}{AUTH-003 --- Log in}
\storymeta{Registered user}{log in with username and password}{I can reach the main application or the pending-approval screen}{\priPzero}{\srcSpec; \srcMock}
\achead
\begin{enumerate}
  \item Given I am on AB, when I enter a valid username and password and tap Enter, and I am approved, then I land on A1 for that community.
  \item Given I enter valid credentials but I am not yet approved, when I tap Enter, then I land on AD (not A1).
  \item Given I enter invalid credentials, when I tap Enter, then I remain on AB and see an error; no session is created.
\end{enumerate}
\authz{Issues an authenticated session on success. Session mechanism is not specified (implementation choice).}
\end{story}

\begin{story}{AUTH-004 --- Forgot password}
\storymeta{Registered user}{use Forgot password on AB}{I can recover access}{\priPone}{\srcMock}
\achead
\begin{enumerate}
  \item Given I am on AB, when I use Forgot password, then a recovery path exists.
  \item The mockup labels the control but does not specify the recovery steps. \textbf{TBD --- Product Decision Required} for exact flow (email reset vs other).
\end{enumerate}
\end{story}

\begin{story}{AUTH-005 --- No Google or social login in V1}
\storymeta{Visitor}{not be required to use Google or other social login}{I can join with username, password, and email only}{\priPzero}{\srcMock}
\achead
\begin{enumerate}
  \item Given I am on AC, when the screen is shown, then it does not offer Google or other third-party login.
  \item The mockup note is: ``we don't need google login etc now.''
\end{enumerate}
\end{story}

\begin{story}{AUTH-006 --- Log out / session end}
\storymeta{Authenticated user}{end my session}{another person using the device cannot continue as me}{\priTBD}{Not specified in S1--S3}
\notes{Login is specified; logout is not drawn or described. Treat presence, placement, and session lifetime as \textbf{TBD --- Product Decision Required}. Do not invent a settings screen for this.}
\end{story}

% ===========================================================================
\section{Epic 2 --- Community Access \& Membership}
% ===========================================================================

\begin{story}{COMM-001 --- Open a community by URL}
\storymeta{Visitor}{open a community at the community's public address}{I land on that community's start page, not another community}{\priPzero}{\srcSpec; \srcTx; \srcMock}
\achead
\begin{enumerate}
  \item Given a community exists with a public address, when I open that address, then I see AA for \emph{that} community (name, picture, Login, Join, scrollable text/images).
  \item Given I open community A's address, when AA renders, then I do not see community B's name, members, or content.
  \item The concrete address shape (path \texttt{domain/community} vs host \texttt{community.domain}) is \textbf{TBD --- COMM-005}.
\end{enumerate}
\end{story}

\begin{story}{COMM-002 --- Community start page content}
\storymeta{Visitor}{see short community information before I log in or join}{I understand which community this is}{\priPzero}{\srcSpec; \srcMock}
\achead
\begin{enumerate}
  \item Given I am on AA, then I see community name, a picture, Login, Join, and further scrollable text and image blocks.
  \item The sources do not specify a CMS for this content. Who edits AA copy is \textbf{TBD} (likely community/platform admin; not drawn).
\end{enumerate}
\end{story}

\begin{story}{COMM-003 --- Community isolation (hard boundary)}
\storymeta{Member of community A}{only see people, groups, chats, and search results that belong to community A}{I am not exposed to communities I never joined}{\priPzero}{\srcTx}
\achead
\begin{enumerate}
  \item Given I am authenticated in community A, when I open chat list, profile search, profile view, or group search, then every result is scoped to community A.
  \item Given I request community B's member list, chats, or profile schema by identifier, when I am not a member of B, then the system denies the request.
  \item The transcript rejects cross-community search: each community is ``totally separate from everything else.''
\end{enumerate}
\authz{Membership in the requested community is required for all member APIs.}
\end{story}

\begin{story}{COMM-004 --- Membership only after manual admin approval}
\storymeta{Community Admin}{decide who enters the community after reviewing the person}{trolls and fake identities are not admitted automatically}{\priPzero}{\srcSpec; \srcTx; \srcMock}
\achead
\begin{enumerate}
  \item Given a user completed Join and is on AD, when they have not been approved, then they are not an approved community member and do not get A1.
  \item Given I am community admin, when I approve the user (via the verification path), then they may use A1 as a member.
  \item Mockup F1 shows a Community checkbox labelled ``Admit to community'' on the ``add this user to groups'' screen. That control, if implemented, must create or confirm community membership.
\end{enumerate}
\deps{VERIFY-001, ADMIN-003, F1/ADMIN-006.}
\end{story}

\begin{story}{COMM-005 --- Path versus subdomain (conflict)}
\storymeta{Platform Admin}{publish a stable public address per community}{members can bookmark and share the community}{\priTBD}{\srcSpec; \srcTx}
\notes{Spec heading says ``sub domain''; written and spoken examples are path-style (\texttt{https://maindomain/community}, \texttt{meetirl.com/bikingstockholm}). \textbf{TBD --- Product Decision Required.} Either choice still requires a first-class ``current community'' on every screen and API.}
\end{story}

\begin{story}{COMM-006 --- Account scope: per community vs platform-wide}
\storymeta{User}{use an account model the product has actually decided}{I am not blocked by an unstated join model}{\priTBD}{Not stated in S1--S3}
\notes{The mockup starts Join from a community URL, which is compatible with either model. The spec never says whether one person joins a second community with the same credentials. \textbf{TBD --- Product Decision Required.} This decision affects unique username scope, AUTH-001 uniqueness, and later multi-community.}
\end{story}

\begin{story}{COMM-007 --- Additional communities later}
\storymeta{Platform Admin}{onboard further communities later, each with its own admin and profile fields}{the platform can host more than one interest circle}{\priPtwo}{\srcSpec; \srcTx}
\achead
\begin{enumerate}
  \item Multiple communities on one platform are specified as later for self-serve onboarding.
  \item The founder stated he will personally set up the first approximately five communities; that is an operations plan, not a V1 self-serve console.
  \item V1 must still keep data shaped \emph{per community} so a second community does not require a redesign.
\end{enumerate}
\end{story}

% ===========================================================================
\section{Epic 3 --- User Profile}
% ===========================================================================
Base fields named in the spec: username (editable); couple/man/woman/other; description (about 4--32k characters); location (country + region list, plus free-text city/village); small chat-icon picture; main profile picture; albums. The mockup adds: three main images; latest login date+time; Score + Notes; Albums; Ads; DM (other profile only); fields that may depend on status.

Do not add age, height, interests, or other fields as hardcoded product fields; those were given as examples of ``many more fields --- later'' or as dynamic schema.

\begin{story}{PROF-001 --- View and edit own profile}
\storymeta{User}{open and edit my own profile}{other members can see who I am}{\priPzero}{\srcSpec; \srcMock}
\achead
\begin{enumerate}
  \item Given I am approved and on A1, when I tap my avatar, then I open my own profile in edit-capable mode.
  \item Given I am on my profile, when I change an editable field and save (save control is implied by ``editable''; exact save UX is not drawn), then the stored values update and later views show the new values.
  \item Mockup labels own profile \textbf{B1} and also labels A1 avatar as \textbf{A3}. Implement one own-profile screen; treat the ID clash as a documentation issue, not two products.
\end{enumerate}
\authz{Only the owner (or a role the sources do not name) may edit. Do not invent admin profile-edit.}
\end{story}

\begin{story}{PROF-002 --- View another member's profile}
\storymeta{User}{open another member's profile from a username, search result, or member list}{I can decide whether to contact them}{\priPzero}{\srcSpec; \srcMock}
\achead
\begin{enumerate}
  \item Given I tap a username in chat, search, member list, or ad (where linked), when the target is in my community, then I see their profile in non-editable mode.
  \item Given I am viewing another profile, then I do not see edit controls, replace-icon, or add-image-to-their-album.
  \item Given the target is not in my community, when I request the profile, then access is denied.
\end{enumerate}
\end{story}

\begin{story}{PROF-003 --- Username}
\storymeta{User}{have an editable username}{others can find and recognise me}{\priPzero}{\srcSpec}
\achead
\begin{enumerate}
  \item Given I am on my profile, when I change username to a value allowed by uniqueness rules, then the new username appears in chats, lists, and search.
  \item Uniqueness scope follows COMM-006 (TBD). Allowed character set is not specified (\textbf{TBD}).
\end{enumerate}
\end{story}

\begin{story}{PROF-004 --- Status: man / woman / couple / other}
\storymeta{User}{set couple, man, woman, or other}{search and profile show my status}{\priPzero}{\srcSpec; \srcMock}
\achead
\begin{enumerate}
  \item Given I edit my profile, when I set status, then I can choose among couple, man, woman, other (checkbox or pulldown --- sources allow either).
  \item Given the mockup note that fields can depend on status, when status changes, then the visible field set may change \emph{if} the community schema defines status-dependent fields. Exact dependency rules are \textbf{TBD} until the first schema is defined.
\end{enumerate}
\end{story}

\begin{story}{PROF-005 --- Long description}
\storymeta{User}{write a long description}{others can read about me}{\priPzero}{\srcSpec}
\achead
\begin{enumerate}
  \item Given I edit my profile, when I enter description text, then the field accepts a long value. The spec states approximately 4--32k characters.
  \item Given I exceed the maximum, when I save, then the system rejects or truncates according to a rule that is not further specified --- implement the stated 32k upper bound unless product revises it.
\end{enumerate}
\end{story}

\begin{story}{PROF-006 --- Location without kilometre distance}
\storymeta{User}{set country, region, and optional city/village}{others can find people in a region}{\priPzero}{\srcSpec; \srcTx}
\achead
\begin{enumerate}
  \item Given I edit location, when I choose country and region, then I pick from lists (Sweden example: on the order of 30 regions). I may also enter free-text city/village.
  \item Given I use profile search, when I filter by region, then results are members of this community in that region. Optional substring match on city/village is described in the transcript.
  \item Kilometre / GPS distance search must not be implemented as V1. The transcript rejected it as too complicated.
\end{enumerate}
\end{story}

\begin{story}{PROF-007 --- Small chat-icon picture}
\storymeta{User}{set a small picture used as my chat avatar}{my posts and chat rows show a recognisable icon}{\priPzero}{\srcSpec; \srcMock}
\achead
\begin{enumerate}
  \item Given I am on my profile, when I replace my icon, then subsequent chat rows and posts use the new icon.
  \item File type, max size, and crop rules are not specified (\textbf{TBD}).
\end{enumerate}
\end{story}

\begin{story}{PROF-008 --- Main profile image(s)}
\storymeta{User}{show main profile image(s)}{others see a primary photo}{\priPzero}{\srcSpec; \srcMock}
\achead
\begin{enumerate}
  \item Spec names a changeable main profile image. Mockup draws \textbf{three} main images, each tappable to a larger view, with verified + date where applicable.
  \item \textbf{TBD --- Product Decision Required:} one main image (spec) vs three (mockup). Do not invent a fourth layout.
  \item Given I tap a main image on another profile, when I am allowed to view it, then I see a larger image.
\end{enumerate}
\end{story}

\begin{story}{PROF-009 --- Latest login shown on profile}
\storymeta{User}{see another member's latest login date and time on their profile}{I have a freshness signal}{\priPone}{\srcMock}
\achead
\begin{enumerate}
  \item The mockup shows ``latest login'' with date+time on the profile screen.
  \item Whether this appears on own profile, other profiles, or both, and whether unapproved users have a login timestamp, is not written in the spec. Treat visibility rules as \textbf{TBD} if the team questions privacy; the mockup does draw it.
\end{enumerate}
\end{story}

\begin{story}{PROF-010 --- Start a DM from another profile}
\storymeta{User}{start a DM from another member's profile}{I can contact them after viewing their information}{\priPzero}{\srcSpec; \srcMock}
\achead
\begin{enumerate}
  \item Given I view another profile (A3), when I use DM / start chat, then I open or create a DM with that user (A4) if I am allowed to contact them.
  \item DM is not shown as a primary action on own profile.
  \item ``Choose what profiles can contact me'' is named in the spec under DM and is not designed. If that rule exists, this action must respect it (SET-002).
\end{enumerate}
\end{story}

\begin{story}{PROF-011 --- Photo albums --- create, rename, list}
\storymeta{User}{add, delete, and rename albums of images}{I can organise pictures on my profile}{\priPzero}{\srcSpec; \srcMock}
\achead
\begin{enumerate}
  \item Given I am on my profile, when I open Albums (C5), then I see my albums as a scrollable list/grid with names.
  \item Given I tap +, when I create an album with a name, then it appears in the list.
  \item Given I long-press an album I own, when I confirm delete, then the album and its images are removed. The mockup requires an ask/confirm.
\end{enumerate}
\end{story}

\begin{story}{PROF-012 --- Album images --- add and delete}
\storymeta{User}{add and delete images inside an album I own}{I control what others may see after unlock}{\priPzero}{\srcSpec; \srcMock}
\achead
\begin{enumerate}
  \item Given I open my album (E1), when I add an image (camera, upload, or other path offered by B4 where applicable), then the image appears with an upload/capture date.
  \item Given I long-press my image, when I confirm delete, then it is removed.
  \item Images taken with the in-app camera show a verified indicator and date (VERIFY-004).
\end{enumerate}
\end{story}

\begin{story}{PROF-013 --- Lock an album with a password}
\storymeta{User}{set a password on an album I own}{only people I choose to tell the password can open it}{\priPzero}{\srcSpec; \srcMock}
\achead
\begin{enumerate}
  \item Given I am in E1, when I set a password, then subsequent views by others require that password (E2).
  \item The owner can open their own album without the other-user password prompt.
  \item How the password is shared (chat, verbally, etc.) is not an in-app grant flow. \textbf{TBD} if a grant/revoke feature is wanted; sources describe typing a password only.
\end{enumerate}
\end{story}

\begin{story}{PROF-014 --- View another user's locked album}
\storymeta{User}{open another user's locked album by entering the password}{I can see pictures they reserved for trusted people}{\priPzero}{\srcSpec; \srcMock}
\achead
\begin{enumerate}
  \item Given I open a locked album (E2), when I have not entered the correct password this viewing session, then I am prompted for the password and images are not shown.
  \item Given I enter the correct password, when the check succeeds, then I see the album images for this viewing session.
  \item Given I enter an incorrect password, when I submit, then images remain hidden.
  \item Failed-attempt lockout is not specified --- do not invent one as a product rule.
\end{enumerate}
\authz{Requester must already be allowed to view that profile (same community). Password is a second gate.}
\end{story}

\begin{story}{PROF-015 --- Dynamic fields render on the profile}
\storymeta{User}{see and fill the fields configured for this community}{my profile matches this community's questions}{\priPzero}{\srcTx; \srcMock}
\achead
\begin{enumerate}
  \item Given community A has a schema of field definitions, when I open a profile in A, then I see those field types (examples drawn: text, checkbox, pulldown, numeric, multi-choice, area) and not community B's unique fields.
  \item Given I own the profile, when I edit a field, then the value is stored against that field definition and this community.
\end{enumerate}
\deps{DYN-001, DYN-002.}
\end{story}

\begin{story}{PROF-016 --- Ads and groups shortcuts on profile}
\storymeta{User}{reach ads and (if admin) add-to-groups from a profile}{I can continue the contact or admin flow}{\priPone}{\srcMock}
\achead
\begin{enumerate}
  \item Mockup profile shows Albums, Ads, Score+Notes, and admin-only ``add to groups'' (F1).
  \item Ads targets: other user's ads C4 vs own ads C9. Whether Ads is V1 is AD-001 (conflict).
  \item F1 is P0 as an admin capability (ADMIN-006) if the control is shown only to admins.
\end{enumerate}
\end{story}

% ===========================================================================
\section{Epic 4 --- Dynamic Profile Fields}
% ===========================================================================
This is an architectural product requirement, not an optional extra. The transcript agrees: do not hardcode the first community's fields; provide a backend so the founder can add a checkbox without a code change. Self-serve UI for other community admins is later.

\begin{story}{DYN-001 --- Field types supported by the backend}
\storymeta{Community / Platform Admin}{define fields by type rather than by a fixed profile table}{I can change the first community without a deploy for each checkbox}{\priPzero}{\srcTx}
\achead
\begin{enumerate}
  \item Given an authorised configurator, when they add a field, then the type is one of the types the sources name: checkbox, text (including long text), image, numeric, pulldown/dropdown, multi-choice, area/location-style, and the search mappings in G1 (multi-select, numeric interval, substring).
  \item Do not treat the first community's eventual labels (e.g.\ ``bike type'') as hardcoded columns in the application.
\end{enumerate}
\end{story}

\begin{story}{DYN-002 --- Schema is per community}
\storymeta{System}{store field definitions per community}{community A and community B can have different questions}{\priPzero}{\srcTx; \srcSpec}
\achead
\begin{enumerate}
  \item Given community A defines field F and community B does not, when a B member opens their profile or search, then F is not shown and is not queryable.
  \item Values are stored per user per field definition (and therefore per community).
\end{enumerate}
\end{story}

\begin{story}{DYN-003 --- Search form is generated from the same schema}
\storymeta{User}{search using controls that match this community's field types}{I know what I can search for}{\priPzero}{\srcTx; \srcMock}
\achead
\begin{enumerate}
  \item Given community A's schema contains a numeric field, when I open G1, then I get a numeric interval (min/max) control for that field --- not a free-text box for the whole profile.
  \item Given a checkbox/multi-select field, when I open G1, then I get a multi-select/checkbox filter.
  \item Given a text field, when I open G1, then I get a substring-in-field control as drawn.
  \item Free-text global search across unknown keywords is rejected in the transcript.
\end{enumerate}
\deps{SEARCH-002.}
\end{story}

\begin{story}{DYN-004 --- Who configures fields in V1}
\storymeta{Platform Admin (founder)}{set up profile fields for the first communities myself}{I can change fields after two weeks without waiting for a deploy}{\priPzero}{\srcTx}
\achead
\begin{enumerate}
  \item V1 requires the \emph{capability} to add/change field definitions without shipping a new hardcoded profile.
  \item The founder stated he will not hand this to other community admins for the first approximately five communities.
  \item A polished self-serve ``community admin field designer'' UI is \textbf{not} required as a V1 user-facing product if an internal/admin path exists for the founder. Exact UI for that path is \textbf{TBD} (no screen is drawn for field configuration).
\end{enumerate}
\end{story}

\begin{story}{DYN-005 --- Self-serve field configuration for other community admins}
\storymeta{Community Admin of a future community}{configure my own profile fields without the founder}{new communities do not need engineering}{\priPtwo}{\srcTx}
\achead
\begin{enumerate}
  \item The transcript places this in the future: ``here's your backend, do it yourself'' will not happen very early.
\end{enumerate}
\end{story}

\begin{story}{DYN-006 --- First community's concrete field list}
\storymeta{Community Admin}{have an agreed field list for community one}{engineering can seed real definitions}{\priTBD}{\srcTx}
\notes{The founder is not 100\% sure what the first profile should contain. Do not invent a biking schema. Seed only fields the product owner supplies.}
\end{story}

% ===========================================================================
\section{Epic 5 --- Profile Search}
% ===========================================================================

\begin{story}{SEARCH-001 --- Open profile search (A5)}
\storymeta{Approved member}{open Find Profiles}{I can discover people I do not already chat with}{\priPzero}{\srcSpec; \srcMock}
\achead
\begin{enumerate}
  \item Given I am approved in this community, when I open Profiles from the shell (mockup bottom tab; transcript: Search icon), then I see A5.
  \item Navigation chrome is NAV-003 (TBD: two vs three tabs). The search \emph{capability} is P0 either way.
\end{enumerate}
\end{story}

\begin{story}{SEARCH-002 --- Filter form G1}
\storymeta{Approved member}{filter profiles with a form built from this community's schema}{I can specify man/woman, age span, region, and other configured fields}{\priPzero}{\srcTx; \srcMock}
\achead
\begin{enumerate}
  \item Given I tap the filter control on A5, when G1 opens, then I see example-type controls: multi-select, numeric interval, substring --- populated from this community's schema, not from another community.
  \item Given I belong only to community A, when G1 renders, then I do not see fields that exist only on community B.
\end{enumerate}
\end{story}

\begin{story}{SEARCH-003 --- Results are community-scoped}
\storymeta{Approved member}{receive only matching members of this community}{I cannot find people outside this circle}{\priPzero}{\srcTx}
\achead
\begin{enumerate}
  \item Given I submit filters, when results return, then every row is an approved member of the current community who matches the filters.
  \item Given no members match, when results return, then the list is empty (empty-state copy is not specified; show an empty list, do not invent members).
  \item Cross-community and free-text global search must not run.
\end{enumerate}
\end{story}

\begin{story}{SEARCH-004 --- Result row content}
\storymeta{Approved member}{see enough on each result to decide whom to open}{I do not have to open every profile}{\priPzero}{\srcMock}
\achead
\begin{enumerate}
  \item Given A5 has results, when a row renders, then it shows username, status, Age1/Age2, district, location, and headline as drawn --- \emph{if} those attributes exist on the schema / profile. If the first community schema omits age, do not hardcode age columns; the mockup is an example layout.
\end{enumerate}
\end{story}

\begin{story}{SEARCH-005 --- Open a profile from a result}
\storymeta{Approved member}{tap a result and open that profile}{I can continue to DM or notes}{\priPzero}{\srcMock}
\achead
\begin{enumerate}
  \item Given I tap a result row, when navigation completes, then I see that user's profile (other-user mode).
\end{enumerate}
\end{story}

\begin{story}{SEARCH-006 --- Persist filters between sessions}
\storymeta{Approved member}{have my profile filters remembered between sessions}{I do not rebuild the form every time}{\priPzero}{\srcMock}
\achead
\begin{enumerate}
  \item Given I set filters on G1, when I leave and later return to profile search in this community on this account, then the last filters are restored. The mockup subtitle is ``save between sessions.''
\end{enumerate}
\end{story}

\begin{story}{SEARCH-007 --- Simple versus advanced search by user level}
\storymeta{User at a given level}{see either a simple or an advanced search}{payment or admin level can gate search depth}{\priTBD}{\srcTx}
\notes{Only example given: simple = age and gender; advanced = more detail. Which level gets which, and whether V1 implements the gate, is \textbf{TBD}. If unimplemented, all approved members see the full community schema form.}
\end{story}

% ===========================================================================
\section{Epic 6 --- Direct Messaging}
% ===========================================================================

\begin{story}{DM-001 --- Chat list with unread counts and sort}
\storymeta{Approved member}{see my DM and group chats with unread counts}{I know where I have new messages}{\priPzero}{\srcSpec; \srcMock}
\achead
\begin{enumerate}
  \item Given I have chats, when A1 loads, then each row shows the chat/group name and unread count when unread $>$ 0.
  \item Given some chats have unread messages, when the list is sorted, then unread chats are above chats with zero unread; zero-unread chats are ordered by most recent activity.
  \item Given I am viewing a chat, when a new message arrives in that chat, then that message does not increase my unread count for that chat.
\end{enumerate}
\end{story}

\begin{story}{DM-002 --- Search among my existing chats}
\storymeta{Approved member}{filter my open chats by substring}{I can find a person among many connections}{\priPzero}{\srcSpec; \srcMock}
\achead
\begin{enumerate}
  \item Given I type in ``search your chats'', when the query is non-empty, then the list shows only my existing chats whose name matches as a substring.
  \item This search must not become community-wide people search (that is A2 / A5).
\end{enumerate}
\end{story}

\begin{story}{DM-003 --- Start or find a DM via plus (A2)}
\storymeta{Approved member}{search a username substring and open a DM}{I can contact someone I met as a name, e.g.\ biker439}{\priPzero}{\srcSpec; \srcMock}
\achead
\begin{enumerate}
  \item Given I tap + (and I am not using the admin-only create-group path), when A2 opens, then I can filter users and groups by substring.
  \item Given I tap a user row, when I continue, then I open a DM (A4) or their profile then DM. Spec: ``click user = choose to start chat or view profile.'' Exact one-step vs two-step is slightly underspecified; both actions must exist. \textbf{P1} on the exact interstitial.
\end{enumerate}
\end{story}

\begin{story}{DM-004 --- Leave, delete history, and block}
\storymeta{User in a DM}{leave the DM and choose whether to delete the chat and/or block the other person}{I can stop an unwanted conversation}{\priPzero}{\srcSpec; \srcMock}
\achead
\begin{enumerate}
  \item Given I am in a DM, when I choose Leave (G2), then I see the options drawn: delete chat + leave + block; leave + block; leave chat.
  \item Given I block the user, when they later try to start a new DM with me, then they cannot.
  \item Given I choose delete chat, when the action completes, then that chat history is removed as described (``the other user cannot go back and look at the messages''). Confirm whether delete is for both sides as the transcript states --- implement as bilateral delete when ``delete chat'' is chosen.
\end{enumerate}
\end{story}

\begin{story}{DM-005 --- Admin: add/remove this DM user on groups I admin}
\storymeta{Group Admin}{from a DM, tick groups I admin to add or remove the other user}{I can admit someone who just messaged ``can I join?'' without opening their profile}{\priPzero}{\srcSpec; \srcTx; \srcMock F1}
\achead
\begin{enumerate}
  \item Given I am a group admin and I am in a DM, when I open add-to-groups (F1), then I see groups I admin with on/off checkboxes, plus a Community ``admit'' checkbox as drawn.
  \item Given I turn a group checkbox on, when saved, then that user becomes a member and the group appears in their A1.
  \item Given I turn a group checkbox off, when saved, then membership is removed and the group disappears from their list.
  \item A non-admin must not see this control.
\end{enumerate}
\end{story}

\begin{story}{DM-006 --- Who may start a DM (``profiles that can contact me'')}
\storymeta{User}{choose what profiles can contact me}{I am not opened to every member by default if the product intends a gate}{\priTBD}{\srcSpec}
\notes{The spec lists this under DM as a bullet and does not define the rule (everyone in community, verified only, level-based, block-list only, etc.). Blocking is specified (DM-004). Additional inbound gates are \textbf{TBD}.}
\end{story}

\newpage
% ===========================================================================
\section{Epic 7 --- Group Chats}
% ===========================================================================

\begin{story}{GROUP-001 --- Community admin creates a group (V1)}
\storymeta{Community Admin (or a user already allowed to create groups)}{create a group with name, info posts, calendar flag, and settings}{members have a room with rules and events}{\priPzero}{\srcSpec later text; \srcTx; \srcMock A9}
\achead
\begin{enumerate}
  \item Given I have create-group rights (A7), when I use the admin control on A1 that opens A9, then I can set name; calendar on/off; Info 1--3 headline+text; new members see history; members see user list; allow post editing; and (drawn) allow direct join.
  \item Given I tap Create, when validation passes, then the group exists in this community, I am group admin, and I may add members via A11/F2 or start empty.
  \item Regular members must not see the create-group entry in V1 if the restrictive rule is confirmed (GROUP-002).
\end{enumerate}
\end{story}

\begin{story}{GROUP-002 --- Regular users cannot create groups in V1}
\storymeta{User}{not create a group unless later allowed}{group creation stays with the community admin}{\priPzero}{\srcSpec ``only community admin\ldots now''; \srcTx}
\achead
\begin{enumerate}
  \item Given I am a normal user, when I open +, then I get A2 (search user/group) and not A9.
  \item \textbf{Conflict:} the spec's opening bullet says users can make a group and be admin. The walkthrough and later spec sentence override that for V1. If product later chooses the opening bullet, this story is revoked.
\end{enumerate}
\end{story}

\begin{story}{GROUP-003 --- Users may create groups later}
\storymeta{User}{create a group when community admin or payment plan allows it}{more rooms can form without the community admin}{\priPtwo}{\srcSpec; \srcTx}
\achead
\begin{enumerate}
  \item Later, community admin can allow others/anyone to create groups, possibly by payment plan / user level. Not V1.
\end{enumerate}
\end{story}

\begin{story}{GROUP-004 --- Add members by username search}
\storymeta{Group Admin}{search a username and add that person to this group}{I can grow the room}{\priPzero}{\srcSpec; \srcMock F2}
\achead
\begin{enumerate}
  \item Given I am group admin, when I open F2, then I can filter users and add them with the add control.
  \item Given I add a user, when the write succeeds, then they see the group on A1.
  \item Avatar/name opens their profile.
\end{enumerate}
\end{story}

\begin{story}{GROUP-005 --- Remove a member}
\storymeta{Group Admin}{remove a user from the group after a confirm}{I can eject someone}{\priPzero}{\srcSpec; \srcMock A11}
\achead
\begin{enumerate}
  \item Given I am on A11, when I use delete on a member and confirm, then membership ends and the group disappears from their A1.
  \item A normal member cannot remove others.
\end{enumerate}
\end{story}

\begin{story}{GROUP-006 --- Member list visibility}
\storymeta{Group member}{see the member list only if the group setting allows it}{the admin can protect the roster}{\priPzero}{\srcSpec; \srcMock}
\achead
\begin{enumerate}
  \item Given ``allow people to see group members list'' is on, when a member opens group info (G3) or the members path, then they see the list.
  \item Given the setting is off, when a non-admin member requests the list, then the API denies it --- hiding the button is not enough.
  \item Group admin can always manage members (A11).
\end{enumerate}
\end{story}

\begin{story}{GROUP-007 --- Custom member status (admin CRM fields)}
\storymeta{Group Admin}{set a short text and three checkboxes on a member (example: paid)}{I can track group-operational status}{\priPzero}{\srcSpec}
\achead
\begin{enumerate}
  \item Given I open a member from the admin list, when I edit status, then I have four fields: one short text and three checkboxes.
  \item These fields are group-membership scoped, not community profile fields.
  \item Who can see these status fields besides the group admin is not specified (\textbf{TBD}; default: admin-only).
\end{enumerate}
\end{story}

\begin{story}{GROUP-008 --- Group info shown to non-members}
\storymeta{User who found a group}{read the group info text before joining}{I know what the room is}{\priPzero}{\srcSpec; \srcMock A2}
\achead
\begin{enumerate}
  \item Given I select a group in A2, when I view group info, then I see the admin-maintained info text (image later).
  \item I do not see the private message history unless I am a member and history rules allow it.
\end{enumerate}
\end{story}

\begin{story}{GROUP-009 --- Rename and delete group}
\storymeta{Group Admin}{rename or delete my group}{I can correct the name or close the room}{\priPzero}{\srcSpec; \srcMock A10}
\achead
\begin{enumerate}
  \item Given I am group admin on A10, when I rename, then the new name appears on A1 and in the chat header.
  \item Given I delete the group, when I confirm, then the group is removed for all members. Confirm copy is not drawn beyond the button --- require a confirm.
\end{enumerate}
\end{story}

\begin{story}{GROUP-010 --- V1 group settings}
\storymeta{Group Admin}{toggle the V1 settings}{each room can behave differently}{\priPzero}{\srcSpec}
\achead
\begin{enumerate}
  \item Given I open group settings, when I save, then these flags persist and take effect immediately for members:
  \begin{itemize}
    \item Allow new members to see history before / when joining.
    \item Allow members to see the member list.
    \item Allow users to edit their own posts.
    \item Show group calendar.
    \item Show sticky posts.
  \end{itemize}
  \item Each flag is enforced server-side, not only in the UI.
\end{enumerate}
\end{story}

\begin{story}{GROUP-011 --- Later group settings}
\storymeta{Group Admin}{use settings that the spec marks later}{future rooms can hide usernames or allow direct join}{\priPtwo}{\srcSpec; \srcTx}
\achead
\begin{enumerate}
  \item Later: allow users to see others' usernames or not.
  \item Later: allow direct join vs message-admin-only. Transcript: V1 they always message the admin. Mockup A9 still draws the checkbox --- \textbf{P1/TBD} whether the control is visible but off, or omitted until later.
\end{enumerate}
\end{story}

\begin{story}{GROUP-012 --- Leave a group}
\storymeta{Group member}{leave the group}{I always have one action as a normal member}{\priPzero}{\srcSpec; \srcTx}
\achead
\begin{enumerate}
  \item Given I am a non-admin member, when I open the chat-name menu, then Leave is available.
  \item Given I leave, when the action completes, then I am not a member and the chat is gone from my A1.
  \item Leave-and-block is specified for DMs, not for groups.
\end{enumerate}
\end{story}

\begin{story}{GROUP-013 --- Join by messaging the group admin (V1 default)}
\storymeta{User}{request access by DMing the group admin}{I can ask to join without a self-serve join}{\priPzero}{\srcSpec; \srcTx; \srcMock A2}
\achead
\begin{enumerate}
  \item Given direct join is not allowed (V1 default), when I select a group in A2, then I may view info and start a DM with the group admin; I do not become a member until an admin adds me.
\end{enumerate}
\end{story}

\begin{story}{GROUP-014 --- History visible to new members when allowed}
\storymeta{New group member}{see prior posts if the admin allowed history}{I can read context (the annoying Wire default is no history)}{\priPzero}{\srcSpec; \srcTx}
\achead
\begin{enumerate}
  \item Given the history setting is on, when I am added, then I can read posts written before I joined.
  \item Given the setting is off, when I am added, then I do not receive pre-join history.
\end{enumerate}
\end{story}

\begin{story}{GROUP-015 --- Co-admins}
\storymeta{Group Admin}{assign co-admins}{others can help run the group}{\priPtwo}{\srcSpec}
\achead
\begin{enumerate}
  \item Assign co-admins is marked later. V1 has a single group admin (the creator / community-appointed admin).
\end{enumerate}
\end{story}

\begin{story}{GROUP-016 --- Maximum group size}
\storymeta{Group Admin}{know whether a size cap exists}{I do not hit an unstated limit}{\priTBD}{\srcSpec}
\notes{The spec writes ``max group size?'' as an open question. Unlimited group \emph{chats} is a platform NFR; per-group headcount is unspecified.}
\end{story}

% ===========================================================================
\section{Epic 8 --- Chat Posts, Reactions \& Comments}
% ===========================================================================
``Post'' is the spec's umbrella word for text, pictures, and comments. DM and group share this behaviour except where a story says group-only.

\begin{story}{POST-001 --- Text post}
\storymeta{Chat member}{send a text post}{others in the chat see it}{\priPzero}{\srcSpec; \srcMock}
\achead
\begin{enumerate}
  \item Given I am a member of the chat, when I send text from the composer, then a post is stored with my user id, date+time, and my small profile picture/icon.
  \item Given I am not a member, when I attempt to post, then the system rejects the write.
  \item Feed order on A4: newest toward the bottom (mockup).
\end{enumerate}
\end{story}

\begin{story}{POST-002 --- Picture post}
\storymeta{Chat member}{post a picture}{others see the image in the feed}{\priPzero}{\srcSpec}
\achead
\begin{enumerate}
  \item Given I add a picture via B4 and send, when the post is created, then members see the image with date+time and my identity.
  \item Video, audio, and collage are later (POST-012).
\end{enumerate}
\end{story}

\begin{story}{POST-003 --- Choose camera, upload, or album image (B4)}
\storymeta{Chat member}{add a picture from camera, new upload, or profile/album}{I can attach a real or existing image}{\priPzero}{\srcMock}
\achead
\begin{enumerate}
  \item Given I tap add-picture, when B4 opens, then I see: Camera (picture will be ``Verified''); Upload new; From profile/album.
  \item Given I choose Camera, when capture succeeds, then the image is stored with in-app-camera = true and a capture timestamp.
  \item Given I upload or pick an album image, when posted, then it is not tagged as in-app-camera (unless that file was originally captured that way --- sources do not specify inheritance; treat a fresh upload as not camera-tagged).
\end{enumerate}
\end{story}

\begin{story}{POST-004 --- React with emoticons}
\storymeta{Chat member}{react to a post with emoticons (not only a single like)}{I can acknowledge a post the way Wire does}{\priPzero}{\srcSpec; \srcMock}
\achead
\begin{enumerate}
  \item Given I tap/long-press a post, when I pick an emoticon, then my reaction is stored and counts update for all members.
  \item Multiple emoticon types are required (Wire-like row).
\end{enumerate}
\end{story}

\begin{story}{POST-005 --- See who reacted}
\storymeta{Chat member}{see who reacted with which emoticon}{reactions are not anonymous counts only}{\priPzero}{\srcSpec; \srcMock B3}
\achead
\begin{enumerate}
  \item Given a post has reactions, when I open the who-reacted / details path from the post menu, then I see users and their emoticons.
\end{enumerate}
\end{story}

\begin{story}{POST-006 --- Comment with original-post snippet}
\storymeta{Chat member}{comment on a post}{my reply is attached to that post}{\priPzero}{\srcSpec; \srcMock}
\achead
\begin{enumerate}
  \item Given I choose Reply/Comment, when I send text, then a comment is stored on that post.
  \item When the comment is shown in the feed, it displays a small, visually dimmed snippet of the original post above or with the comment.
\end{enumerate}
\end{story}

\begin{story}{POST-007 --- Delete own post}
\storymeta{Chat member}{delete my own post and its comments}{I can remove something I should not have sent}{\priPzero}{\srcSpec}
\achead
\begin{enumerate}
  \item Given I own the post, when I delete it, then the post and its comments (and reactions) are removed for everyone.
  \item Own-delete is always allowed; it is not gated by the edit setting.
\end{enumerate}
\end{story}

\begin{story}{POST-008 --- Edit own post if the group allows it}
\storymeta{Group member}{edit my own post only when the group setting allows editing}{I can fix typos when the admin permits it}{\priPzero}{\srcSpec}
\achead
\begin{enumerate}
  \item Given I am in a group, the allow-edit setting is on, and I own the post, when I use Edit, then I can change the post content.
  \item Given the setting is off, when I am a non-admin owner, then Edit is not offered and the API rejects edit.
  \item DM edit: the edit setting is listed under group settings. Whether DMs are always editable is \textbf{TBD}. Do not silently copy the group flag onto DMs.
\end{enumerate}
\end{story}

\begin{story}{POST-009 --- Group admin deletes any post}
\storymeta{Group Admin}{delete another member's post}{I can moderate the room}{\priPzero}{\srcSpec; \srcTx}
\achead
\begin{enumerate}
  \item Given I am group admin of this group, when I long-press any post, then Delete is available for others' posts as well as mine.
  \item This is not gated by a setting. Admins can always delete posts.
  \item There is no admin-delete concept in a DM.
\end{enumerate}
\end{story}

\begin{story}{POST-010 --- Long-press / post menu (B3)}
\storymeta{Chat member}{open a large post menu}{power actions stay hidden until needed}{\priPzero}{\srcSpec; \srcTx; \srcMock}
\achead
\begin{enumerate}
  \item Given I long-press (or use the spec's click-hold) a post, when the menu opens, then it is a large/full-width sheet, not a tiny popover.
  \item Menu items from the mockup: emoticons; who reacted; Details; Edit (own or admin); Delete (own or admin-others); Reply; Copy; Cancel.
  \item Copy and Details are Wire leftovers also drawn on the mockup --- include them as drawn. Details content beyond who-reacted is not specified (\textbf{TBD} if Details $\neq$ who-reacted).
\end{enumerate}
\end{story}

\begin{story}{POST-011 --- Open poster profile from username}
\storymeta{Chat member}{tap a username on a post and open that profile}{I can move from chat to profile}{\priPzero}{\srcSpec; \srcMock}
\achead
\begin{enumerate}
  \item Given I tap the poster's name or icon, when they are in my community, then I open their profile.
\end{enumerate}
\end{story}

\begin{story}{POST-012 --- Video, audio, collage later}
\storymeta{Chat member}{post video, audio, or collage}{richer media}{\priPtwo}{\srcSpec; \srcTx}
\end{story}

\begin{story}{POST-013 --- Temporary / limited-view posts later}
\storymeta{Chat member}{send a post limited by time or number of views}{Snapchat-like ephemerality}{\priPtwo}{\srcSpec; \srcTx}
\end{story}

\begin{story}{POST-014 --- Filter group posts by profile parameters later}
\storymeta{Group member}{see only posts from people matching profile filters}{I can skim an Ads-like room}{\priPtwo}{\srcSpec; \srcTx}
\notes{This is the transcript's V1-ads concept (``just a post + later filter''). Dedicated ad screens are a separate conflict (Epic 15).}
\end{story}

% ===========================================================================
\section{Epic 9 --- Verification}
% ===========================================================================

\begin{story}{VERIFY-001 --- Manual admin verification of new users}
\storymeta{Community Admin}{manually decide whether a new user is who they say they are}{admission is not algorithmic}{\priPzero}{\srcSpec; \srcTx; \srcMock AD}
\achead
\begin{enumerate}
  \item Given a user is on AD, when they have filled required profile information and contacted the admin, then I can review and approve or not approve.
  \item Given I have not approved them, when they log in, then they remain on AD / unapproved (AUTH-003).
  \item There is no automated identity vendor in V1.
\end{enumerate}
\end{story}

\begin{story}{VERIFY-002 --- Unapproved user fills profile and DMs the admin}
\storymeta{Unapproved user}{fill my profile and message the community admin}{the admin can review me}{\priPzero}{\srcMock AD}
\achead
\begin{enumerate}
  \item Given I am on AD, then I see copy that I must fill the profile, message the admin for verification, and that I need to fill all info to be approved.
  \item Given I tap Create/Edit profile, when it opens, then I reach own-profile edit.
  \item Given I tap Chat with admin, when it opens, then I have a DM with the community admin (A4).
\end{enumerate}
\end{story}

\begin{story}{VERIFY-003 --- In-app camera capture for verification evidence}
\storymeta{User}{take a live picture in the app (ID, card, or handwritten note with date and name)}{the admin can see a capture that is harder to AI-upload}{\priPzero}{\srcTx; \srcSpec}
\achead
\begin{enumerate}
  \item Given I use Camera for this evidence, when the capture UI opens, then I cannot satisfy this step by picking an arbitrary gallery file in that same action (frontend forces capture).
  \item Transcript examples: handwritten note with today's date and name; holding ID/card in a live snapshot. Spec image rule also says date. Require date; name-on-note is from the transcript --- include both in helper copy unless product shortens it.
  \item The captured file is stored with in-app-camera = true and timestamp.
\end{enumerate}
\end{story}

\begin{story}{VERIFY-004 --- Verified indicator on every image}
\storymeta{Any viewer}{see whether an image was taken with the internal app camera, plus date}{I can treat camera-tagged images as a stronger trust signal}{\priPzero}{\srcSpec}
\achead
\begin{enumerate}
  \item Given an image was captured via the in-app camera, when it is shown anywhere (chat, profile, album, ad), then the UI shows that it was taken by the internal camera and shows the date.
  \item Given an image was a normal upload, when it is shown, then it does not claim to be camera-verified.
\end{enumerate}
\end{story}

\begin{story}{VERIFY-005 --- Honest meaning of ``verified picture''}
\storymeta{System / UX copy}{not claim cryptographic proof of identity}{we do not over-promise}{\priPzero}{\srcTx (implicit); product integrity}
\achead
\begin{enumerate}
  \item UI copy may say the picture was taken with the in-app camera and show the date.
  \item UI copy must not claim the image is impossible to fake. Sources present this as an added layer / deterrent for admin review.
\end{enumerate}
\end{story}

\begin{story}{VERIFY-006 --- AI-assisted verification later}
\storymeta{Community Admin}{use AI later (e.g.\ face-picture check)}{some review can be assisted}{\priPtwo}{\srcTx}
\end{story}

% ===========================================================================
\section{Epic 10 --- Sticky Posts \& Calendar}
% ===========================================================================

\begin{story}{STICKY-001 --- View sticky information in a group}
\storymeta{Group member}{open named info items (rules, FAQ, etc.) from the group chat header}{I can read room rules without hunting the history}{\priPzero}{\srcSpec; \srcTx; \srcMock}
\achead
\begin{enumerate}
  \item Given I am in a group (A4B) and ``show sticky posts'' is on, when the header shows Info items (examples: Rules, FAQ; admin-chosen headlines), then I can open one and read the long text on B2.
  \item Given the setting is off, when I am a normal member, then sticky controls are not shown and the info API is not available to me.
  \item DMs do not have official sticky posts.
\end{enumerate}
\end{story}

\begin{story}{STICKY-002 --- Admin create, edit, delete sticky posts}
\storymeta{Group Admin}{add, edit, and delete sticky posts with a headline and long body}{new members can be pointed at rules}{\priPzero}{\srcSpec; \srcMock A9}
\achead
\begin{enumerate}
  \item Given I am group admin, when I edit Info 1--3 (headline + text), then members see the updated headlines and bodies.
  \item Spec says ``info 1+2'' / multiple allowed; mockup draws three. Support at least two; three matches the later mockup. Do not invent an unbounded CMS beyond what is drawn unless product expands it.
\end{enumerate}
\end{story}

\begin{story}{CAL-001 --- View group calendar}
\storymeta{Group member}{open the group calendar and see events on a month view}{I can see that there is a ride on Tuesday}{\priPzero}{\srcSpec; \srcMock B1}
\achead
\begin{enumerate}
  \item Given I am a member and ``show group calendar'' is on, when I open the calendar, then I see a month (choose month/year), dates with events marked, and tapping a date lists that day's events (header + text).
  \item There is no community-wide or platform-wide calendar in the sources.
  \item DMs have no calendar.
\end{enumerate}
\end{story}

\begin{story}{CAL-002 --- Admin create, edit, delete calendar events}
\storymeta{Group Admin}{add, edit, and delete calendar items}{I can publish group events}{\priPzero}{\srcSpec; \srcMock G4}
\achead
\begin{enumerate}
  \item Given I am group admin, when I use + or tap an event, then G4 lets me set date, headline, and text, and Save or Delete.
  \item Non-admins cannot create/edit/delete events. The sources do not allow members to propose events --- do not invent that.
  \item Recurring events, reminders, and notifications tied to events are not specified.
\end{enumerate}
\end{story}

% ===========================================================================
\section{Epic 11 --- Personal Notes \& Ratings}
% ===========================================================================
Two note kinds exist. Do not conflate them. Reputation and profile score are later and are not these private stars.

\begin{story}{NOTE-001 --- Personal sticky note on a chat}
\storymeta{User}{write a long private note on a chat that only I can see}{I can remember ``Sven is annoying'' while I chat}{\priPzero}{\srcSpec; \srcTx; \srcMock B5}
\achead
\begin{enumerate}
  \item Given I am in a chat, when I open the fixed private-notes control, then I can add/edit/delete a long text note scoped to that chat and to me.
  \item Given any other member opens the same chat, when they look, then they cannot see my note or that it exists.
\end{enumerate}
\end{story}

\begin{story}{NOTE-002 --- Private note on another profile}
\storymeta{User}{write a note on another profile that only I can see}{I can keep a mini-CRM}{\priPzero}{\srcSpec; \srcMock}
\achead
\begin{enumerate}
  \item Given I view another profile, when I save a note, then it is stored as author $\rightarrow$ subject and shown only to me.
  \item The profile owner must never see the note text or an indicator that a note exists.
\end{enumerate}
\end{story}

\begin{story}{NOTE-003 --- Private 1--5 star rating}
\storymeta{User}{set a 1--5 star rating on another profile that only I can see}{I can rank people for myself}{\priPzero}{\srcSpec; \srcMock}
\achead
\begin{enumerate}
  \item Given I view another profile, when I set stars 1--5, then only I can see that rating.
  \item The rating is not aggregated, averaged, or shown to the owner in V1.
  \item Mockup B5 draws notes and stars on one screen titled with group name or profile name --- one UI may serve chat-note and profile-note+rating; privacy rules still differ by target.
\end{enumerate}
\end{story}

\begin{story}{NOTE-004 --- Profile completeness score later}
\storymeta{User}{receive an automatic score for filling fields and verified pictures}{complete profiles are rewarded later}{\priPtwo}{\srcTx}
\notes{Examples given only as thinking: 200 characters = points; main picture = points; verified main picture = more points. Not in the current specification to build.}
\end{story}

\begin{story}{NOTE-005 --- Reputation later}
\storymeta{User}{be scored by others' opinions of my behaviour}{mass DMs and bad actors can be limited later}{\priPtwo}{\srcTx}
\notes{Could later stop someone from contacting people. Not specified for V1.}
\end{story}

\newpage
% ===========================================================================
\section{Epic 12 --- Navigation \& Main Application UI}
% ===========================================================================

\begin{story}{NAV-001 --- Main screen A1 as home after approval}
\storymeta{Approved member}{land on the chat list as the main application view}{chat is the home base}{\priPzero}{\srcSpec; \srcMock}
\achead
\begin{enumerate}
  \item Given I log in as approved, when the app opens, then A1 shows: avatar, title Chats, + (and admin create control if A7), search-your-chats, chat rows, bottom navigation as decided in NAV-003.
\end{enumerate}
\end{story}

\begin{story}{NAV-002 --- Plus opens join-chat search (full-screen)}
\storymeta{User}{open a large + sheet to find a user or group}{starting a conversation is easy to tap}{\priPzero}{\srcTx; \srcMock A2}
\achead
\begin{enumerate}
  \item Given I tap +, when A2 opens, then it uses a large portion of the screen (transcript: do not use a tiny dropdown).
  \item Entries depend on create-group rights (GROUP-001/002).
\end{enumerate}
\end{story}

\begin{story}{NAV-003 --- Bottom navigation structure (conflict)}
\storymeta{User}{use the bottom navigation the product actually chose}{I am not trained on a layout that will be reversed}{\priTBD}{\srcTx vs \srcMock}
\achead
\begin{enumerate}
  \item \textbf{Alternative A (transcript):} two items --- Chats (this list) and Search (profile search form). Own profile = top-left avatar. Founder said the third icon is not needed.
  \item \textbf{Alternative B (mockup 260922):} three items --- Chats, Profiles (A5), Ads (A6). Own profile still = avatar.
  \item Engineering must not pick A or B silently. \textbf{TBD --- Product Decision Required.}
\end{enumerate}
\end{story}

\begin{story}{NAV-004 --- Wire-like, low-chrome interaction}
\storymeta{User}{see few icons and discover extra actions via long-press}{the app feels limited but easy}{\priPzero}{\srcTx}
\achead
\begin{enumerate}
  \item Given a normal member uses a chat, when they do not long-press, then they can still post and react without opening admin menus.
  \item QA should flag screens that add icon bars not present in Wire or the handwritten map without a source.
\end{enumerate}
\end{story}

\begin{story}{NAV-005 --- Screen inventory (mockup IDs)}
\storymeta{Designer / frontend}{implement the drawn screens rather than inventing new primary pages}{the flow matches Peter's map}{\priPzero}{\srcMock}
\notes{See Appendix~A. Blank or unfinished: G3 is almost empty (name, headline, member list). Do not invent G3's remaining content. Ads screens exist on paper but conflict with the transcript (Epic 15).}
\end{story}

% ===========================================================================
\section{Epic 13 --- Settings}
% ===========================================================================
Only settings named by the sources. Do not add notification, theme, or language settings as V1 product requirements.

\begin{story}{SET-001 --- Group settings screen}
\storymeta{Group Admin}{open group settings as a list of on/off options}{I can configure the room}{\priPzero}{\srcSpec; \srcTx; S4 as pattern only}
\achead
\begin{enumerate}
  \item Given I am group admin, when I open settings, then I see the V1 flags in GROUP-010 as explicit toggles (transcript: radio/toggle style like Wire options).
  \item Do not copy Wire's iOS notification, calls, or protocol settings as product settings.
  \item Later flags (GROUP-011) are omitted or disabled unless product explicitly includes the A9 ``direct join'' checkbox as visible-but-later.
\end{enumerate}
\end{story}

\begin{story}{SET-002 --- Inbound DM contact rules}
\storymeta{User}{choose what profiles can contact me}{unwanted DMs can be reduced}{\priTBD}{\srcSpec}
\notes{Named under Chat / DM in the spec. No UI, no rule set. Blocking (DM-004) is the only specified inbound control. Do not invent a preference panel.}
\end{story}

\begin{story}{SET-003 --- No additional user settings invented}
\storymeta{User}{only see settings the sources name}{the app stays minimal}{\priPzero}{Methodology}
\achead
\begin{enumerate}
  \item V1 must not add theme, language, notification, privacy-policy, or ``who can see my profile'' screens unless product later specifies them. Profile visibility inside a community is implied: members can view other members' profiles.
\end{enumerate}
\end{story}

% ===========================================================================
\section{Epic 14 --- Notifications}
% ===========================================================================

\begin{story}{NOTIF-001 --- Browser notifications later}
\storymeta{User}{receive a notification when a new message arrives}{I do not miss DMs}{\priPtwo}{\srcSpec; \srcTx}
\achead
\begin{enumerate}
  \item Spec: turn on/off iPhone + Android notifications when a new message arrives --- marked later.
  \item Transcript: possible on the web, but the notification is from the browser (e.g.\ Safari), showing the website name, not a native app notification.
  \item V1 has no notification requirement.
\end{enumerate}
\end{story}

% ===========================================================================
\section{Epic 15 --- Advertising}
% ===========================================================================

\begin{story}{AD-001 --- V1 ads concept (conflict)}
\storymeta{Product Owner}{decide what Ads means in V1}{engineering does not build two products}{\priTBD}{\srcTx vs \srcMock}
\achead
\begin{enumerate}
  \item \textbf{Alternative A (transcript):} no specific ad functionality. An ad is a normal post in a group (likely named Ads). Filtering that feed by profile parameters is later (POST-014). Monetisation is later.
  \item \textbf{Alternative B (mockup 260922):} dedicated screens --- A6 all ads; C4 one user's ads; C9 my ads; D1 sort/filter (remembered between sessions); D2 publish/edit own ad; D3 view others' ad (not editable) with Profile and DM.
  \item Do not implement B as if A did not exist. \textbf{TBD --- Product Decision Required.}
\end{enumerate}
\end{story}

\begin{story}{AD-002 --- Dedicated ad list and publish (only if Alternative B is chosen)}
\storymeta{User}{browse ads and publish my own}{I can post ``looking for Saturday woods riders in Stockholm''}{\priPone}{\srcMock}
\achead
\begin{enumerate}
  \item Given Alternative B is approved, when I open A6, then I see ad rows: headline, username, status, ages, district, location, published date.
  \item Given I publish via D2, when I save, then I set Looking for (profile-type pulldown), headline, text, district (Sweden-area pulldown), location (free text), image (via B4). Age1, Age2, Date, Status are shown and \textbf{not editable} on the ad (from profile).
  \item Given I view another's ad (D3), when I use Profile or DM, then I reach that profile or a DM.
  \item Given I own the ad (D2), when I publish changes or delete, then the list updates.
  \item Filter D1 remembers values between sessions. Field types drawn: multi-select, numeric interval, substring, status, looking for, district, location.
\end{enumerate}
\notes{Priority is P1 because the screens are drawn but contradicted by the transcript.}
\end{story}

\begin{story}{AD-003 --- Monetised advertising later}
\storymeta{Platform / Community Admin}{charge for ads later}{ads become a revenue feature}{\priPtwo}{\srcTx}
\achead
\begin{enumerate}
  \item Transcript: keep ads as a monetised feature for the future. Do not build billing into V1 ads even if Alternative B ships as a content type.
\end{enumerate}
\end{story}

% ===========================================================================
\section{Epic 16 --- Payments \& Monetisation}
% ===========================================================================

\begin{story}{PAY-001 --- Payment methods later}
\storymeta{Payer (any role)}{pay by credit card and Swedish Swish}{local and card payment exist later}{\priPtwo}{\srcSpec; \srcTx}
\end{story}

\begin{story}{PAY-002 --- Multi-level fees later}
\storymeta{Platform / Community / Group Admin / User}{pay along the hierarchy}{community admin pays platform (base + per user or \%); group admins may pay community; users may pay to use a community and/or a group}{\priPtwo}{\srcSpec; \srcTx}
\end{story}

\begin{story}{PAY-003 --- Community admin as payment receiver later}
\storymeta{Community Admin}{receive a profit share later}{running a community can be compensated}{\priPtwo}{\srcTx}
\notes{Founder has not figured out how. Do not design a split engine.}
\end{story}

\begin{story}{PAY-004 --- User levels gated by payment later}
\storymeta{User}{gain rights (e.g.\ advanced search, create groups) via payment level}{monetisation can unlock features}{\priPtwo}{\srcSpec; \srcTx}
\notes{Depends on the TBD level matrix. Not V1.}
\end{story}

% ===========================================================================
\section{Epic 17 --- Administration}
% ===========================================================================

\begin{story}{ADMIN-001 --- Platform Admin responsibilities}
\storymeta{Platform Admin}{own the platform and onboard communities}{there is a single top operator (today: Petter)}{\priPzero}{\srcSpec; \srcTx}
\achead
\begin{enumerate}
  \item Given I am platform admin, then I can have a community created and can act as community admin for the first community.
  \item A full platform console (billing, all-community dashboard) is not specified. Do not invent one.
\end{enumerate}
\end{story}

\begin{story}{ADMIN-002 --- Community Admin configures this community's profile schema}
\storymeta{Platform Admin / Community Admin}{define this community's profile fields}{members see the right questions}{\priPzero}{\srcTx}
\achead
\begin{enumerate}
  \item Follow DYN-004: V1 needs a way for the founder to add/change field definitions for a community. No dedicated screen is drawn.
  \item Self-serve for other community admins is DYN-005 (P2).
\end{enumerate}
\end{story}

\begin{story}{ADMIN-003 --- Community Admin verifies and admits users}
\storymeta{Community Admin}{review AD users and admit them to the community}{membership stays manual}{\priPzero}{\srcTx; \srcMock}
\achead
\begin{enumerate}
  \item Given an unapproved user DMs me and has a profile (and camera evidence if requested), when I admit them (F1 community checkbox or equivalent), then they become an approved member and reach A1 on next login.
  \item Rejection / request-more-info UX is not drawn (\textbf{TBD} for exact decline path; not admitting them is enough for V1).
\end{enumerate}
\end{story}

\begin{story}{ADMIN-004 --- Community Admin creates groups in V1}
\storymeta{Community Admin}{create groups}{rooms exist only when I allow them}{\priPzero}{\srcSpec; \srcTx}
\deps{GROUP-001, GROUP-002.}
\end{story}

\begin{story}{ADMIN-005 --- Group Admin surface (chat-name menu)}
\storymeta{Group Admin}{reach admin actions from the chat name}{admin stays inside the chat, like Wire}{\priPzero}{\srcSpec; \srcTx}
\achead
\begin{enumerate}
  \item Given I tap the group name, when I am admin, then I can: edit sticky posts; rename; edit group info text; delete group; add users; list members (click to profile, DM, or edit custom status); remove people; add/edit/delete calendar items; open group settings.
  \item Co-admins later (GROUP-015).
\end{enumerate}
\end{story}

\begin{story}{ADMIN-006 --- Admit to community and assign groups (F1)}
\storymeta{Community or Group Admin}{from a person, tick community admission and groups}{chat remains the admin interface}{\priPzero}{\srcMock; \srcTx}
\achead
\begin{enumerate}
  \item Given I open F1 for a user, when I toggle Community, then I admit or (if off) do not admit them to the community. Whether unchecking removes an already-approved member is not specified (\textbf{TBD}).
  \item Group checkboxes follow DM-005.
\end{enumerate}
\end{story}

\begin{story}{ADMIN-007 --- Group Admin always deletes posts}
\storymeta{Group Admin}{delete any post in my group}{I can moderate}{\priPzero}{\srcSpec}
\deps{POST-009.}
\end{story}

% ===========================================================================
\section{Epic 18 --- Security, Privacy \& Data Access}
% ===========================================================================
Sources say ``High security'' and ``GDPR compliant'' and ``European servers'' without naming mechanisms. Stories below translate those statements and the isolation/privacy rules. They do not mandate JWT vs session, specific password policy, rate limits, or audit logs.

\begin{story}{SEC-001 --- Authorise every write by role and scope}
\storymeta{System}{reject actions the caller is not allowed to perform}{login alone is not enough}{\priPzero}{\srcSpec roles; \srcTx}
\achead
\begin{enumerate}
  \item Given a valid session, when the caller deletes a post they do not own and they are not group admin of that group, then the write is rejected.
  \item Given a valid session, when the caller lists members of a group whose member-list setting is off and they are not admin of that group, then the read is rejected.
  \item Given a valid session for community A, when the caller queries community B data, then the request is rejected.
\end{enumerate}
\end{story}

\begin{story}{SEC-002 --- Community isolation on all member data}
\storymeta{Member}{never receive another community's profiles, messages, or schema}{isolation bugs are treated as privacy defects}{\priPzero}{\srcTx}
\deps{COMM-003.}
\end{story}

\begin{story}{SEC-003 --- Private notes and ratings never leak to the subject}
\storymeta{User}{have my notes and stars about someone stay mine}{the mini-CRM is one-directional}{\priPzero}{\srcSpec}
\achead
\begin{enumerate}
  \item Given User A stored a note or rating about User B, when User B loads their own profile or any API about themselves, then that note/rating is absent.
  \item Given User C is not the author, when C requests A's note about B, then access is denied.
\end{enumerate}
\end{story}

\begin{story}{SEC-004 --- Locked album password is a second gate}
\storymeta{User}{keep album images hidden until the correct album password is entered}{trust is shared manually}{\priPzero}{\srcSpec}
\deps{PROF-013, PROF-014.}
\end{story}

\begin{story}{SEC-005 --- Blocked users cannot start a new DM}
\storymeta{User}{remain unreachable by someone I blocked}{leave+block works}{\priPzero}{\srcSpec}
\deps{DM-004.}
\end{story}

\begin{story}{SEC-006 --- European servers and GDPR}
\storymeta{User / operator}{have the service hosted on European servers and treated as GDPR-compliant}{the spec's legal constraints are met}{\priPzero}{\srcSpec; \srcTx}
\achead
\begin{enumerate}
  \item Hosting must be in Europe as stated.
  \item GDPR-compliant is stated; export, deletion, lawful basis, DPA, and retention are \textbf{not} enumerated. Those mechanics are \textbf{TBD} as compliance design, not as invented user stories.
\end{enumerate}
\end{story}

\begin{story}{SEC-007 --- High security without a mandated mechanism}
\storymeta{Operator}{treat security as a stated requirement}{auth and storage are not casual}{\priPzero}{\srcSpec}
\notes{No password policy, 2FA, rate limit, or audit log is specified. Implementation may apply standard practice; those extras are not V1 product stories unless product adds them.}
\end{story}

\begin{story}{SEC-008 --- Image file constraints}
\storymeta{User}{upload images the system will accept}{albums and posts have a defined limit}{\priTBD}{Not specified}
\notes{Image storage approach and max file sizes are not addressed in the sources.}
\end{story}

\newpage
% ===========================================================================
\section{End-to-End User Flows}
% ===========================================================================
Flows use the mockup IDs. A diamond marks a TBD branch.

\subsection{Flow 1 --- New user $\rightarrow$ community $\rightarrow$ profile}
\begin{verbatim}
Visitor
  -> Opens community URL
     -> System shows AA (name, picture, Login, Join)
Visitor
  -> Taps Join
     -> AC: username, password, password, email
  -> Taps Send
     -> Account created (unapproved)
     -> Email with login link -> AB
User
  -> Logs in (not approved)
     -> AD
  -> Create/Edit profile -> own profile fields (schema of this community)
  -> Chat with admin -> A4 DM
Community Admin
  -> Reviews profile / camera evidence
  -> Admits via F1 Community checkbox (or equivalent)
     -> Membership = approved
User
  -> Next login -> A1
\end{verbatim}
Business writes: Account; ProfileFieldValues; CommunityMembership (unapproved$\rightarrow$approved); optional Verification image (camera-tagged).

\subsection{Flow 2 --- User $\rightarrow$ profile search $\rightarrow$ profile $\rightarrow$ DM}
\begin{verbatim}
Approved member on A1
  -> Opens Profiles / Search  [NAV-003: tab vs icon]
     -> A5 result list (previous G1 filters restored)
  -> Opens G1, sets community-schema filters, returns
     -> Query scoped to this community only
  -> Taps a row
     -> Other profile (non-editable)
  -> Optional: private note / stars (only author can read)
  -> Taps DM
     -> A4 DM created or opened
     -> If inbound-contact rule exists (SET-002 TBD), gate applies
\end{verbatim}

\subsection{Flow 3 --- Existing chats $\rightarrow$ open $\rightarrow$ send message}
\begin{verbatim}
Approved member on A1
  -> Optionally substring-filters own chats
  -> Taps a row (unread badge if any)
     -> A4 / A4B
  -> Types text or adds picture (B4)
  -> Sends
     -> AuthN + AuthZ (member?)
     -> Post stored with timestamp + avatar
     -> Other members' unread counts increment
     -> Transport (websocket / poll / SSE) = implementation choice
\end{verbatim}

\subsection{Flow 4 --- Search group $\rightarrow$ info $\rightarrow$ join request}
\begin{verbatim}
Approved member
  -> Taps + -> A2
  -> Substring filter
  -> Taps a group
     -> Group info text (public)
  -> If direct join allowed [V1 default: NO; mockup checkbox TBD]
        -> Member immediately
     Else
        -> DM group admin (A4)
Group Admin
  -> Uses F1 / F2 / DM shortcut
     -> GroupMembership created
     -> Group appears on requester's A1
     -> History visible only if GROUP-014 setting is on
\end{verbatim}

\subsection{Flow 5 --- Group admin manages members}
\begin{verbatim}
Group Admin in group A4B
  -> Taps chat name -> admin menu / A10
  -> User list A11
  -> Tap avatar -> profile
  -> Tap delete -> confirm -> membership removed
  -> Tap + -> F2 search -> add
  -> Tap member -> custom status (1 text + 3 checkboxes)
Alternative (faster):
  Admin in a DM -> F1 tick/untick groups (+ community admit)
\end{verbatim}

\subsection{Flow 6 --- Group admin sticky post}
\begin{verbatim}
Group Admin
  -> A9/A10 or chat-name menu
  -> Edit Info N headline + long text
     -> StickyPost upsert
Member in A4B (setting show sticky = on)
  -> Taps headline (e.g. Rules)
     -> B2 full text
\end{verbatim}

\subsection{Flow 7 --- Group admin calendar event}
\begin{verbatim}
Group Admin
  -> Calendar from A4B (or A9 calendar flag)
     -> Month view B1
  -> Taps + or an event -> G4
  -> Date, headline, text -> Save or Delete
Member (show calendar = on)
  -> Opens calendar
  -> Month with marked dates
  -> Taps a date -> list of events
\end{verbatim}

\subsection{Flow 8 --- User verification}
\begin{verbatim}
Unapproved user on AD
  -> Completes required profile fields
  -> Opens in-app camera (cannot satisfy with a random gallery pick)
  -> Holds note (date + name) and/or ID as requested
  -> Image stored: camera-tagged + timestamp
  -> DMs community admin
Community Admin
  -> Reviews profile + tagged image
  -> Approves (manual)
     -> User may use A1
Copy must not claim the photo is unfakeable.
\end{verbatim}

\subsection{Flow 9 --- Locked album}
\begin{verbatim}
Owner on own profile
  -> Albums C5 -> E1
  -> Sets password; adds images (verified flag if camera)
Viewer on other profile
  -> Albums C8 -> E2
  -> Password prompt
  -> Correct -> images this session
  -> Incorrect -> still locked
No in-app ``grant access'' flow is specified.
\end{verbatim}

\subsection{Flow 10 --- Admin configures community profile fields}
\begin{verbatim}
Platform Admin (founder), first communities
  -> Uses the field-definition capability (UI not drawn --- TBD)
  -> Adds types: checkbox, text, numeric, image, pulldown, multi-choice, ...
     -> ProfileFieldDef rows for this community only
Member
  -> Own profile shows those fields (PROF-015)
  -> G1 search form shows matching filter widgets (DYN-003)
Other community admins
  -> Self-serve designer = later (DYN-005)
\end{verbatim}

% ===========================================================================
\section{Non-Functional Requirements}
% ===========================================================================

\subsection{Explicit in the sources}
\begin{tabularx}{\textwidth}{@{}l l Y@{}}
\toprule
\textbf{ID} & \textbf{Requirement} & \textbf{Source / notes} \\
\midrule
NFR-001 & Web-based, mobile-first & Spec, transcript \\
NFR-002 & Extremely easy / fast (Wire-like; a slow WordPress-like platform was rejected) & Spec, transcript \\
NFR-003 & European servers & Spec \\
NFR-004 & GDPR compliant & Spec (mechanics TBD) \\
NFR-005 & High security & Spec (mechanism not specified) \\
NFR-006 & Handle 100 concurrent users now & Spec \\
NFR-007 & Handle thousands of users within 2--3 years & Spec \\
NFR-008 & Unlimited group chats and DM chats & Spec \\
NFR-009 & Community isolation as a correctness/privacy property & Transcript \\
\bottomrule
\end{tabularx}

\subsection{Assignment / inferred --- not source product requirements}
\begin{itemize}
  \item \textbf{PWA / service worker / offline:} ``PWA'' appears as a project label, not as an offline or installability requirement in S1--S3. Treat installability and offline as \textbf{TBD} unless the assignment separately mandates them.
  \item \textbf{Realtime transport:} not specified. WebSockets, polling, or SSE are implementation choices.
  \item \textbf{Next.js / Supabase / TypeScript:} implementation context, not NFRs from the client.
  \item \textbf{HTTPS / hashed passwords:} reasonable engineering practice under ``high security,'' not a named product story.
\end{itemize}

% ===========================================================================
\section{V1 Scope}
% ===========================================================================

\begin{longtable}{@{}p{2.1cm}p{4.6cm}p{3.6cm}p{1.6cm}p{2.2cm}@{}}
\toprule
\textbf{ID} & \textbf{Feature} & \textbf{Stories} & \textbf{Pri} & \textbf{Status} \\
\midrule
\endhead
V1-01 & Join, email link, login, pending AD & AUTH-001--005, VERIFY-002 & P0 & In scope \\
V1-02 & Community URL + isolation + manual admit & COMM-001--004, ADMIN-003 & P0 & In scope \\
V1-03 & Own/other profile, base fields, albums, lock & PROF-001--007, 010--015 & P0 & In scope \\
V1-04 & Dynamic field backend + search mapping & DYN-001--004 & P0 & In scope \\
V1-05 & Form-based community profile search & SEARCH-001--006 & P0 & In scope \\
V1-06 & Chat list, unread sort, in-list search, + & DM-001--003, NAV-001--002 & P0 & In scope \\
V1-07 & DM leave / delete / block & DM-004 & P0 & In scope \\
V1-08 & Admin add-to-groups from person/DM & DM-005, ADMIN-006 & P0 & In scope \\
V1-09 & Admin-only group create, members, settings & GROUP-001--002, 004--010, 012--014 & P0 & In scope \\
V1-10 & Posts, images, camera tag, react, comment, edit/delete & POST-001--011 & P0 & In scope \\
V1-11 & Manual verification + camera tag everywhere & VERIFY-001--005 & P0 & In scope \\
V1-12 & Sticky posts + group calendar & STICKY-001--002, CAL-001--002 & P0 & In scope \\
V1-13 & Private chat notes; profile notes; 1--5 stars & NOTE-001--003 & P0 & In scope \\
V1-14 & Minimal Wire-like chrome & NAV-004--005 & P0 & In scope \\
V1-15 & EU hosting, GDPR statement, high-level security & SEC-001--007, NFR-001--009 & P0 & In scope \\
V1-16 & Main-image count (1 vs 3) & PROF-008 & TBD & Blocked on decision \\
V1-17 & Latest login on profile & PROF-009 & P1 & Mockup only \\
V1-18 & Bottom nav 2 vs 3 & NAV-003 & TBD & Blocked \\
V1-19 & Ads dedicated UI & AD-002 & P1 & Blocked on AD-001 \\
\bottomrule
\end{longtable}

% ===========================================================================
\section{Future / Later}
% ===========================================================================
\begin{longtable}{@{}p{2.3cm}p{5.4cm}p{6.4cm}@{}}
\toprule
\textbf{ID} & \textbf{Feature} & \textbf{Stories} \\
\midrule
\endhead
F-01 & Video / audio / collage posts & POST-012 \\
F-02 & Temporary / limited-view posts & POST-013 \\
F-03 & Filter group posts by profile parameters & POST-014 \\
F-04 & Browser / OS notifications & NOTIF-001 \\
F-05 & Users create groups; payment-gated creation & GROUP-003, PAY-004 \\
F-06 & Direct join (if not taken from mockup) & GROUP-011 \\
F-07 & Hide other usernames & GROUP-011 \\
F-08 & Co-admins & GROUP-015 \\
F-09 & Self-serve field UI for other community admins & DYN-005 \\
F-10 & Self-serve multi-community onboarding & COMM-007 \\
F-11 & Reputation score & NOTE-005 \\
F-12 & Profile completeness score & NOTE-004 \\
F-13 & Monetised ads & AD-003 \\
F-14 & Payments: card, Swish, multi-level, profit share & PAY-001--003 \\
F-15 & AI-assisted verification & VERIFY-006 \\
F-16 & Google / social login & AUTH-005 (explicitly not now) \\
\bottomrule
\end{longtable}

% ===========================================================================
\section{TBD / Product Decisions Required}
% ===========================================================================
Only questions the sources actually open or contradict.

\begin{enumerate}[leftmargin=1.6em]
  \item \textbf{NAV-003} Bottom nav: 2 tabs (transcript) vs 3 tabs (mockup).
  \item \textbf{COMM-006} Account created per community vs one platform account that joins communities.
  \item \textbf{COMM-005} Path URL vs true subdomain.
  \item \textbf{GROUP-002} Confirm admin-only group creation (recommended reading of later spec + transcript).
  \item \textbf{DYN-006} Concrete field list for community one.
  \item \textbf{SEARCH-007 / user levels} Permission matrix for levels 1--3 besides simple vs advanced search.
  \item \textbf{AD-001} Ads as normal group posts vs dedicated A6/C4/C9/D1--D3.
  \item \textbf{GROUP-011 / A9} Show ``allow direct join'' in V1 or omit until later.
  \item \textbf{PROF-008} One main image vs three.
  \item \textbf{SET-002} Meaning of ``choose what profiles can contact me.''
  \item \textbf{GROUP-016} Max group size (spec asks).
  \item \textbf{SEC-008} Image storage and max sizes.
  \item \textbf{PROF-013} Whether album password sharing stays fully out-of-band.
  \item \textbf{AUTH-004 / AUTH-006} Forgot-password steps; logout/session lifetime.
  \item \textbf{POST-008} Are DMs editable, given the edit flag is a group setting?
  \item \textbf{GROUP-007} Who besides the group admin sees custom member status.
  \item \textbf{ADMIN-006} Does unchecking Community on F1 revoke membership?
  \item \textbf{SEC-006} GDPR export/delete/retention mechanics.
  \item \textbf{AUTH-002} Login-link token lifetime.
  \item \textbf{Facebook-style ``I agree'' / join questions} Liked in the meeting, \textbf{not} written into the spec --- out of V1 unless product adds them.
  \item \textbf{PWA offline/install} Assignment label vs unspecified in S1--S3.
  \item \textbf{Mockup ID collisions} B1, A3, B5 --- documentation only; one screen each.
\end{enumerate}

% ===========================================================================
\section{Requirements Traceability}
% ===========================================================================
Every meaningful source statement should appear here. ``Unconverted'' means a story cannot be precise until a decision.

\begin{longtable}{@{}p{5.8cm}p{2.0cm}p{3.4cm}p{2.4cm}@{}}
\toprule
\textbf{Source requirement} & \textbf{Epic} & \textbf{Story} & \textbf{When} \\
\midrule
\endhead
Login / create user & E1 & AUTH-001, AUTH-003 & V1 \\
Join fields + email login link; no Google now & E1 & AUTH-001, 002, 005 & V1 \\
Forgot password label & E1 & AUTH-004 & TBD flow \\
Unapproved users go to AD not A1 & E1, E9 & AUTH-003, VERIFY-002 & V1 \\
Community at maindomain/community (wording clash) & E2 & COMM-001, COMM-005 & V1 + TBD \\
Community start page info + login/join & E2 & COMM-002 & V1 \\
Each community totally separate; no cross search & E2, E5 & COMM-003, SEARCH-003 & V1 \\
Manual selection who enters & E2, E9, E17 & COMM-004, VERIFY-001, ADMIN-003 & V1 \\
Multiple communities later; founder sets first $\sim$5 & E2, E4 & COMM-007, DYN-004, DYN-005 & V1 capability / later UI \\
Account per community vs platform & E2 & COMM-006 & TBD \\
Username editable & E3 & PROF-003 & V1 \\
Couple/man/woman/other & E3 & PROF-004 & V1 \\
Description 4--32k & E3 & PROF-005 & V1 \\
Location country/region + city; no km & E3 & PROF-006 & V1 \\
Small chat icon; main image & E3 & PROF-007, PROF-008 & V1 / TBD count \\
Albums add/delete/rename; lock password & E3 & PROF-011--014 & V1 \\
Every image: in-app camera + date & E8, E9 & POST-003, VERIFY-004 & V1 \\
Dynamic field types (cb, text, image, numeric, \ldots) & E4 & DYN-001--003 & V1 \\
Admin customizes fields for the community & E4 & DYN-004, ADMIN-002 & V1 (founder path) \\
Search = fill a profile form; not free text & E5 & SEARCH-002, DYN-003 & V1 \\
Search only inside community & E5 & SEARCH-003 & V1 \\
A5 rows + G1 persist filters & E5 & SEARCH-004--006 & V1 \\
Simple vs advanced search by level & E5 & SEARCH-007 & TBD \\
User levels 1--3 & E2 & SEARCH-007, PAY-004 & TBD / later \\
Chat list, unread top, substring search & E6, E12 & DM-001, DM-002, NAV-001 & V1 \\
+ search username/group substring & E6, E7 & DM-003, GROUP-013 & V1 \\
Click user: start chat or view profile & E6 & DM-003, PROF-002 & V1 \\
Leave DM: delete and/or block & E6 & DM-004, SEC-005 & V1 \\
Admin add user to groups from DM & E6, E17 & DM-005, ADMIN-006 & V1 \\
Choose what profiles can contact me & E13 & SET-002, DM-006 & TBD \\
Only community admin creates groups now & E7 & GROUP-001, GROUP-002 & V1 \\
Users can make groups (opening bullet) & E7 & GROUP-002 conflict & Conflict \\
Users create groups later / payment & E7 & GROUP-003 & Later \\
Add/delete users; member list protect & E7 & GROUP-004--006 & V1 \\
Custom status 1 text + 3 checkboxes & E7 & GROUP-007 & V1 \\
Group info text for non-members & E7 & GROUP-008 & V1 \\
Rename/delete group & E7 & GROUP-009 & V1 \\
Group settings V1 list & E7, E13 & GROUP-010, SET-001 & V1 \\
Username visibility; direct join later & E7 & GROUP-011 & Later / TBD UI \\
Leave group is the always-on member action & E7 & GROUP-012 & V1 \\
History for new members setting & E7 & GROUP-014 & V1 \\
Co-admins later & E7 & GROUP-015 & Later \\
Max group size? & E7 & GROUP-016 & TBD \\
Text post + date/time + user pic & E8 & POST-001 & V1 \\
Picture post; camera post verified & E8 & POST-002, POST-003 & V1 \\
React with emoticons; list who reacted & E8 & POST-004, POST-005 & V1 \\
Comment + original snippet & E8 & POST-006 & V1 \\
Delete own (+ comments); edit if allowed & E8 & POST-007, POST-008 & V1 \\
Admin delete others' posts & E8 & POST-009 & V1 \\
Click username $\rightarrow$ profile & E8 & POST-011 & V1 \\
Video/audio/collage later & E8 & POST-012 & Later \\
Limited-view posts later & E8 & POST-013 & Later \\
Filter posts by profile params later & E8 & POST-014 & Later \\
Personal sticky note in chat & E11 & NOTE-001 & V1 \\
Note on profile; 1--5 stars private & E11 & NOTE-002, NOTE-003 & V1 \\
Reputation later; profile score later & E11 & NOTE-004, NOTE-005 & Later \\
In-app camera; handwritten note/ID; tag image & E9 & VERIFY-003--005 & V1 \\
Manual admin verification & E9 & VERIFY-001 & V1 \\
AI identification later & E9 & VERIFY-006 & Later \\
Sticky posts headline + long text & E10 & STICKY-001, STICKY-002 & V1 \\
Group calendar; admin CRUD events & E10 & CAL-001, CAL-002 & V1 \\
Bottom nav 2 vs 3 & E12 & NAV-003 & TBD \\
Avatar $\rightarrow$ own profile; + full screen & E12 & NAV-001, NAV-002 & V1 \\
Apple not Microsoft; long-press & E12 & NAV-004, POST-010 & V1 \\
Notifications later; browser-level & E14 & NOTIF-001 & Later \\
Ads = group posts (transcript) vs A6--D3 (mockup) & E15 & AD-001, AD-002 & TBD \\
Monetised ads later & E15 & AD-003 & Later \\
Payments card + Swish; multi-level; profit share & E16 & PAY-001--003 & Later \\
Platform / community / group admin roles & E17 & ADMIN-001--007 & V1 \\
High security; GDPR; EU servers & E18 & SEC-006, SEC-007, NFR-003--005 & V1 + TBD mechanics \\
100 concurrent now; thousands in 2--3 years & NFR & NFR-006, NFR-007 & V1 / growth \\
Unlimited chats & NFR & NFR-008 & V1 \\
Web, mobile-first, very fast & NFR & NFR-001, NFR-002 & V1 \\
PWA offline/install & NFR & assignment only & TBD \\
Facebook ``I agree'' / join quiz & --- & meeting aside & Not in spec \\
\end{longtable}

\subsection{Source items that cannot yet become precise stories}
\begin{itemize}
  \item Exact first-community schema (founder unsure).
  \item User-level rights beyond one search example.
  \item GDPR operational rights (export/erasure).
  \item Image size/storage.
  \item Logout and password-reset steps.
  \item Ads V1 shape.
  \item Bottom navigation.
  \item Account and URL models.
\end{itemize}

% ===========================================================================
\section{QA / Acceptance Considerations}
% ===========================================================================
QA should derive cases from Given/When/Then on every P0 story. The following are cross-cutting tests that are easy to miss.

\begin{enumerate}[leftmargin=1.6em]
  \item \textbf{Isolation:} create two communities with different schemas. A member of A must never see B's users, chats, schema, or search hits --- including by guessing IDs.
  \item \textbf{Pending user:} login after join must open AD, not A1, until admit.
  \item \textbf{Search generation:} add a numeric field to community A only; G1 in A shows a min/max control; G1 in B does not.
  \item \textbf{Member list setting:} with the setting off, a non-admin API call for members is 403/denied, not an empty list pretending to be success if the data exists.
  \item \textbf{History setting:} add a user after posts exist; they see history only if the flag is on.
  \item \textbf{Edit setting:} owner edit succeeds only when the group flag is on; admin delete always succeeds.
  \item \textbf{Private CRM:} user B's session never contains A's note or stars about B (UI and API).
  \item \textbf{Album gate:} wrong password never returns image bytes.
  \item \textbf{Camera tag:} upload path must not display ``taken with in-app camera''; capture path must.
  \item \textbf{Block:} blocked user cannot open a new DM.
  \item \textbf{Create group:} normal user + menu has no Create; community admin does.
  \item \textbf{Unread:} viewing a chat does not increment unread for the viewer; other members do increment.
  \item \textbf{Verified copy:} no screen claims an image cannot be faked.
\end{enumerate}

Do not write tests for P2 stories as V1 exit criteria.

% ===========================================================================
\section{Implementation-Readiness Summary}
% ===========================================================================
The product is ready to implement as a \textbf{community-scoped, mobile-first chat + profiles} system if the team treats community isolation, dynamic fields, and manual verification as load-bearing, and if it does not start work on payments, reputation, notifications, or social login.

\subsection{Safe to build now}
Authentication and AD; community-scoped data; profiles (base + dynamic values); albums + password gate; A1 chat list; DM and group messaging; posts/reactions/comments; admin moderation; sticky posts; calendar; V1 group settings; form search; private notes/stars; in-app camera metadata; F1/F2 member admin.

\subsection{Do not build as if decided}
Three-tab shell vs two-tab shell; dedicated ads product; user-created groups; direct join; hardcoded first-community columns; cross-community anything; kilometre search; native push; Swish/card; reputation/profile score.

\subsection{Recommended implementation order (from the team's own method)}
The transcript asked for user stories and frontend/mockups first, then backend. Suggested sequence after this document is approved:
\begin{enumerate}
  \item Community-scoped routing + shell (A1) with seed data.
  \item Auth + AD + admit.
  \item Profile screens + field renderer.
  \item Groups + membership + settings.
  \item Posts / comments / reactions.
  \item Sticky + calendar.
  \item Search form generated from schema.
  \item Camera tag + verification review.
  \item Everything in Future / Later --- not in the first sprint set.
\end{enumerate}

\subsection{Sign-off}
This document is internally consistent with S1--S3 as of 22 September 2026. It is a draft until the product owner resolves the TBD list, especially NAV-003, COMM-005, COMM-006, AD-001, DYN-006, and GROUP-002.

\newpage
\appendix
\section{Handwritten screen inventory (Mockup 260922)}
\begin{longtable}{@{}p{1.6cm}p{4.4cm}p{8.2cm}@{}}
\toprule
\textbf{ID} & \textbf{Screen} & \textbf{Notes} \\
\midrule
\endhead
--- & Legend & Square = if status/role; circle = go to screen \\
AA & Community start & URL, name, picture, Login, Join, scrollable about \\
AB & Login & Username, password, Enter, Forgot password $\rightarrow$ A1 or AD \\
AC & Join & Username, 2$\times$ password, email, Send $\rightarrow$ email link AB \\
AD & New user pending & Fill profile (B1/own); DM admin (A4) \\
A1 & Main chats & Avatar, +, search chats, unread sort, bottom nav \\
A7 & Admin on A1 & Extra create-group control $\rightarrow$ A9 \\
A2 & Join a chat & User/group substring; DM or request/join \\
A9 & Create group & Name, calendar, info 1--3, V1 settings, Create \\
A10 & View/edit group & Same form + Delete group \\
A11 & Group users admin & Profile, delete (ask), + add \\
F2 & Add users to group & Filter, add, open profile \\
F1 & Add this user to groups & Community admit + group checkboxes \\
A4 & Chat (DM) & Feed, composer, leave, notes; no calendar/sticky \\
A4B & Chat (group) & + calendar, info, sticky \\
B3 & Post menu & Reactions, details, edit, delete, reply, copy \\
G3 & Group info & Unfinished: name, headline, member list \\
G2 & Leave DM & Delete+block / leave+block / leave \\
B2 & Info text & Headline + long body \\
B5 & Private notes & Chat or profile; text + stars + Save \\
B4 & Add picture & Camera verified / upload / album \\
B1 & Group calendar & Month, marked dates, event list \\
G4 & Edit event & Date, headline, text, save, delete \\
A5 & Find profiles & Result cards, filter icon \\
G1 & Filter profiles & Persist between sessions; typed controls \\
A3/B1 & Profile & Other vs own (ID clash); 3 images; dynamic fields \\
C5/C8 & Album list & Own vs other; + and long-press delete on own \\
E1/E2 & Album view & E2 password first; verified+date on images \\
A6/C4/C9 & Ads lists & All / one user / mine --- conflict with transcript \\
D2/D3 & Publish / view ad & Editable vs not; profile fields locked on ad \\
D1 & Filter ads & Remember between sessions \\
\bottomrule
\end{longtable}

\section{Wire screenshots --- visual language only}
The specification's Wire photos are \textbf{not} additional features. Copy density, avatars, unread badges, search-in-list, tappable title, full-width action sheet, add-participant search, and toggle-list pattern. Do not implement Wire calls, self-deleting messages, protocol details, or iOS notification settings.

\section{Dependency map (logical)}
\begin{verbatim}
COMM-001 current community
   -> AUTH-001 Join / AUTH-003 Login
   -> COMM-004 membership
        -> NAV-001 A1
        -> PROF-* / DYN-* values
        -> SEARCH-*  (needs DYN schema)
        -> DM-* / GROUP-* / POST-*
             -> STICKY-* / CAL-* (group only)
        -> NOTE-*
        -> VERIFY-* (also sits on AD before membership)

GROUP-001 create
   -> GROUP-004 membership
        -> POST-*
        -> GROUP-010 settings gate history, list, edit, calendar, sticky

DYN-001 field types
   -> DYN-002 per community
        -> PROF-015 render
        -> DYN-003 / SEARCH-002 form
\end{verbatim}

\section{Glossary}
\begin{tabularx}{\textwidth}{@{}l Y@{}}
\toprule
\textbf{Term} & \textbf{Meaning in this product} \\
\midrule
Community & Isolated interest circle with its own members, groups, and profile schema \\
Group & Chat room inside one community \\
Post & Any feed item: text, image, or comment \\
Verified image & Tagged as captured with the in-app camera + date; not a legal identity proof \\
Sticky post & Admin-authored info item (headline + long text) on a group \\
Private note & Author-only text, either on a chat or on a profile \\
Private rating & Author-only 1--5 stars on a profile; not reputation \\
Reputation & Later, others' judgement of behaviour \\
Profile score & Later, automatic completeness points \\
A4B & Group chat view (calendar + official info) \\
\bottomrule
\end{tabularx}

\end{document}
